\documentclass[11pt,a4paper]{article}

\usepackage[utf8]{inputenc}
\usepackage[T1]{fontenc}
\usepackage{lmodern}
\usepackage[margin=2.5cm]{geometry}
\usepackage{amsmath,amssymb,amsthm,mathtools}
\usepackage{booktabs,array,longtable,tabularx}
\usepackage{enumitem}
\usepackage{microtype}
\usepackage{xcolor}
\usepackage[
  colorlinks=true,
  linkcolor=blue!45!black,
  citecolor=blue!45!black,
  urlcolor=blue!45!black,
  pdftitle={Scoped O.1 on the Cone, v6:
  Executed Apex Records, a Closed Constant-Profile Slavnov-Taylor Row,
  and Ward-Identity Boundaries},
  pdfauthor={Ngo Cheung}
]{hyperref}

\newcommand{\RR}{\mathbb R}
\newcommand{\CC}{\mathbb C}
\newcommand{\ZZ}{\mathbb Z}
\newcommand{\NN}{\mathbb N}
\newcommand{\HH}{\mathcal H}
\newcommand{\ka}{\kappa}
\newcommand{\dd}{\mathrm d}
\newcommand{\ii}{\mathrm i}
\newcommand{\id}{\operatorname{id}}
\newcommand{\Reg}{M_{\mathrm{reg}}}
\newcommand{\WW}{\mathcal W}
\newcommand{\Fr}{\mathrm F}
\newcommand{\CR}{\mathrm{CR}}
\newcommand{\Ccb}{\mathcal C}
\newcommand{\Hreg}{(\mathrm H_{\rm reg})}
\newcommand{\HI}{(\mathrm H_I)}
\newcommand{\Span}{(\mathrm{Span})}
\newcommand{\Jcond}{(\mathrm J)}
\newcommand{\eps}{\varepsilon}
\newcommand{\abs}[1]{\lvert #1\rvert}
\newcommand{\norm}[1]{\lVert #1\rVert}
\newcommand{\pair}[2]{\langle #1,#2\rangle}
\newcommand{\ket}[1]{\lvert #1\rangle}
\newcommand{\bra}[1]{\langle #1\rvert}
\newcommand{\proj}[1]{\lvert #1\rangle\langle #1\rvert}
\newcommand{\loc}[1]{\nolinkurl{#1}}
\newcommand{\NOTCLOSED}{\texttt{NOT\_CLOSED}}
\newcommand{\OPEN}{\textsf{OPEN}}
\newcommand{\SUPPLIED}{\textsf{SUPPLIED}}
\newcommand{\CONDITIONAL}{\textsf{CONDITIONAL}}
\newcommand{\UNVALIDATED}{\textsf{UNVALIDATED}}
\newtheorem*{evidence}{Evidence}

\DeclareMathOperator{\Tr}{Tr}
\DeclareMathOperator{\supp}{supp}
\DeclareMathOperator{\diag}{diag}
\DeclareMathOperator{\im}{im}
\DeclareMathOperator{\rank}{rank}
\DeclareMathOperator{\TV}{TV}
\DeclareMathOperator{\ResUV}{Res_{\!UV}}
\newcommand{\ApTr}{\operatorname{Tr}_{\rm ap}}

\newcolumntype{L}[1]{>{\raggedright\arraybackslash}p{#1}}
\newcolumntype{Y}{>{\raggedright\arraybackslash}X}

\theoremstyle{plain}
\newtheorem{theorem}{Theorem}[section]
\newtheorem{proposition}[theorem]{Proposition}
\newtheorem{lemma}[theorem]{Lemma}
\newtheorem{corollary}[theorem]{Corollary}

\theoremstyle{definition}
\newtheorem{definition}[theorem]{Definition}
\newtheorem{assumption}[theorem]{Assumption}
\newtheorem{inputstatement}[theorem]{Supplied input}
\newtheorem{example}[theorem]{Example}

\theoremstyle{remark}
\newtheorem{remark}[theorem]{Remark}

\numberwithin{equation}{section}
\setlist{itemsep=3pt,topsep=5pt}
\setcounter{tocdepth}{2}
\emergencystretch=2em
\allowdisplaybreaks

\newcommand{\scopebox}[2]{%
\par\medskip
\noindent\fbox{\parbox{\dimexpr\linewidth-2\fboxsep-2\fboxrule\relax}{%
\textbf{#1}\par\smallskip #2}}
\par\medskip
}

\newcommand{\separation}{%
\par\smallskip\noindent
\emph{Separation boundary.}
The cone, extension, regulator, and external-field class are prescribed.
A free inverse, operator limit, one-point finite part, finite instrument,
or documentary comparison is not an interacting two-point Ward theorem.
No Einstein dynamics, physical tensor Hessian, spacetime emergence,
or inherited quantum admission follows.
}

\title{\textbf{Scoped O.1 on the Cone}\\[5pt]
\Large Executed Apex Records, a Closed Constant-Profile
Slavnov--Taylor Row, Weighted Apex Coefficients,
and Ward-Identity Boundaries\\[6pt]
\large An artifact-verified revision of the v5 proof-level audit}
\author{Ngo Cheung}
\date{18 October 2026\\
\small Version 6; revision of \loc{XVI-O1-CONE-v5.tex}}

\begin{document}
\maketitle

\begin{abstract}
Version~6 is an artifact-verified revision of the v5 audit of the
constant-transverse-profile calculation for the scoped O.1 problem
on \(C_\alpha\times\RR^2\).
It keeps v5's structure, its evidence classes, and its conditional
theorems. What it changes is that the executable record behind v4
is now archived, rerun, and audited, and it adds the following.

\smallskip\noindent\textbf{New in v6.}
\begin{enumerate}[label=\textup{(\arabic*)},leftmargin=1.6em]
\item \emph{Catalogue re-derived.}
Every entry of the v5 coefficient catalogue agrees exactly with the
executed spectral reductions of the frozen XVI-P vertex arrays.
The executed record also contains three things v4 never printed:
the \((12)(12)\) entries
\(\widehat B_{\mathrm{D01}}=\widehat B_{\mathrm{D03}}
=-8c_1P^2\,\ka^2/(16\pi)\); the complete D02 tadpole array
\(\tfrac{\ka^2}{2\pi}(c_0+2c_1+c_2)\) times a fixed tensor;
and \(\widehat B_{\mathrm{D12}}=0\).
Mixed-helicity monomials occur only on vector-type keys, which are
excluded from the record, not set to zero.
The D13 transverse source trace is computed to be zero.

\item \emph{The \(hh\) Slavnov--Taylor failure is diagnosed.}
The v4 failure comes from the coordinate weight of the antifield
\(K'_{(03)}\). The correct weight is \(-2\); v4 scanned only
positive weights.
With conventions derived from the frozen vertex, the
constant-profile \(hh/Kc\) row closes on CR and on F.
The closure is exact as a rational identity in \(\alpha\), with
\(\lambda=-2\ii\). This value is fixed in advance by the tree-level
data, not fitted. UNIT closes trivially.

\item \emph{A second apex coefficient.}
We prove the \(r^2\)-weighted apex coefficient
\[
\mathcal D_b=-\frac{11\alpha}{180}+\frac{B_2(b)}{3\alpha}
-\frac{B_4(b)}{6\alpha^3},
\qquad
\ApTr\bigl[r^2f(A_b)\bigr]=-\mathcal D_bf'(0).
\]
Its singular-channel shift is \(2\nu(1-\nu^2)/3\).
It passes exact orbifold and twisted-orbifold checks, and it has no
flat CR remnant.

\item \emph{An exact obstruction for \(\HI\) in D01, and its
resolution.}
The \(S=2\to3\) composite trace has the non-integrable diagonal
\(r^{-2\nu-2}\), so \(\HI\) fails as an \(L^2\) statement.
Its unique homogeneous (Mellin) continuation, however, vanishes
identically in \(\nu\). This is shown in closed form and confirmed
independently by quadrature.

\item \emph{Invariant structure of the \(hh\) poles.}
The constant-profile \(hh\) poles decompose into six
\(SO(1,1)\)-invariant structures.
In the complete \(hh\) sum the purely longitudinal structures
cancel, and the gauge variance enters only through
\(b_6+b_7=-\ka^2V/(4\pi)\).
The CR--F difference is gauge invariant:
\(-\nu\ka^2P^2f_0f_0'/(4\pi)\).
\end{enumerate}
A point-by-point table answers every issue v5 raised against v4.
An appendix lists full SHA-256 digests and rerun commands.

\smallskip\noindent\textbf{Still open.}
The 26 Ward combinations and the 34-dimensional projection;
the six exact values;
the D04 quadratic-profile part of the second \(\phi\phi\) value;
the evanescent coefficients;
the 18-dimensional engine map;
an exact \(\Span\) certificate;
and the vector-type keys.
None of these is reported as zero.
The scope remains one loop at \(O(\ka^2)\), the massless two-point
sector, frozen de Donder gauge, and the XVI-P prescription.
The cone and the CR branch remain inputs, and the
destructive-readout analogy supplies no premise.
\(\texttt{bulk\_symmetry\_restored}=\texttt{false}\), and
O.1 on the cone remains \(\NOTCLOSED\).
\end{abstract}


\tableofcontents

%======================================================================
\section{Question, scope, and evidence classes}
\label{sec:introduction}
%======================================================================

\subsection{The scoped question}

The supplied v4 sheet studies transport of the separately scoped
XVI-Q record Q11 to a prescribed conical product, with:
\begin{enumerate}[label=\textup{(\roman*)}]
\item one-loop order and coupling order \(O(\ka^2)\);
\item a two-point sector and vanishing declared scalar mass;
\item frozen de Donder gauge;
\item dimensional reduction with quasi-\(d\) loop momenta and
quasi-four-dimensional vertex indices;
\item the frozen XVI-P prescription and the p22 contract;
\item curvature-regular (CR) and componentwise Friedrichs (F)
apex realizations compared without selecting either by readout.
\end{enumerate}
XVI-P, XVI-Q, and XVI-R are frozen comparison records
\cite{CheungP,CheungQ,CheungR}.
The later supplied p23 excerpt is not silently substituted for the
p22 contract.

The principal issue is not whether a conical heat coefficient exists.
It is whether the complete differentiated graph, on its declared domain
and in its declared prescription, reduces to that coefficient and obeys
the required interacting identities.

\subsection{What is proved, retained, or conditional}

This revision uses four evidence classes.
v6 adds the last one.
\begin{description}[style=nextline,leftmargin=0pt]
\item[Locally proved.]
The normalized scalar apex coefficient; the singular-channel heat
shift; the new \(r^2\)-weighted coefficient and its shift; a
conditional graph-reduction theorem; the longitudinal UV
normalization; algebraic identities between the printed polynomials;
and the stated access/flux distinctions.

\item[Executed and artifact-verified in v6.]
Statements carrying the tag \textsf{[exec]}.
Each was produced by a program in \loc{code/} and logged in
\loc{logs/v4/} or \loc{logs/v6/}, and was rerun or checked in this
revision.
Its full SHA-256 digest is listed in Appendix~\ref{app:artifacts}.
The vertex inputs are read-only copies of the frozen XVI-P/Q
files; their digests equal the originals.
Exact statements are exact rational or symbolic identities in
\(\alpha\) or \(\nu\).
Numerical statements give their precision.
An executed result is a reproducible computation, not an
independent review.

\item[Supplied computational record.]
v4 reports that v6 did not rerun, together with the historical
registry intake.
These are kept as reports.

\item[Conditional graph identification.]
Equality of a printed polynomial with a particular frozen graph
integral still needs the graph-specific hypotheses of
Assumption~\ref{ass:reduction}.
v6 discharges several of these for the constant-profile record;
the ledger in Section~\ref{sec:ledger} lists them.
No proof is reconstructed from a routine name.
\end{description}

\scopebox{Submission and proof boundary (v6)}{%
v6 now archives the vertex arrays, routing, prefactors, and logs
that v5 found missing.
The constant-profile catalogue is therefore an executed result on
the archived frozen inputs, and no longer just a supplied report.
It is still not an independently refereed computation.
The spectral reduction of each graph remains conditional on
Assumption~\ref{ass:reduction}.
Only (R1) and parts of (R3) are discharged here
(Sections~\ref{sec:catalogue}, \ref{sec:HI}).
(R2), (R4), and (R5) are still the stated hypotheses.}

The coefficient values printed in v4 and v5 are unchanged.
v6 adds entries but alters none.

\subsection{What v6 adds}

\begin{enumerate}[label=\textup{(\alph*)}]
\item \emph{\(hh\) Slavnov--Taylor row}
(Section~\ref{sec:ST}).
Derived conventions, an \emph{a priori} constant \(\lambda\), and
exact closure on CR and F for the constant profile.
The v4 failure is diagnosed, and UNIT closes.
\item \emph{Higher apex coefficient}
(Section~\ref{sec:weighted}).
The \(r^2\)-weighted coefficient \(\mathcal D_b\), with checks.
The quadratic-profile \(\phi\phi\) value is partial: the D05a part
is \(0\) and the D04 part is open.
\item \emph{\((\mathrm H_{I,\mathrm{D01}})\)}
(Section~\ref{sec:HI}).
The exact \(L^2\) obstruction, and an exact vanishing Mellin value.
\item \emph{Evanescent coefficients}
(Section~\ref{sec:missing}).
A pole-level statement only; the coefficients themselves are open.
\item \emph{\(hh\) map}
(Section~\ref{sec:structures}).
Six invariant structures, the complete D02/D12 arrays, and the
coefficient vectors.
The 18-dimensional engine map is open.
\item \emph{\(\Span\) certificate}
(Section~\ref{sec:missing}).
Open, with the precise reason.
\item \emph{Response table}
(Section~\ref{sec:response}).
Every v5 issue, answered point by point.
\end{enumerate}

\subsection{Programme non-transfer}

Connections to I-B, I-B-G, I-B-G-S, I-B-G-S-B, I-B-VP, and I-C are
separations, not additional analytical premises or a declared path.
They do not discharge O.1, open the downstream XV-B admission exit,
or resolve the unrelated H2/I3/J2 characteristic-product obligations.

The historical H2 ending in XIV-H is retained version-specifically;
the narrower I2 internal-amplitude construction in XIV-I is not
withdrawn. The complete joint maps at I3/J2 and their subsequent
admission chain remain separate
\cite{CheungCharacteristic,CheungXVB}.
Missing values are never set to zero.

%======================================================================
\section{Prescribed cone, fields, and branch domains}
\label{sec:conventions}
%======================================================================

\subsection{Geometry and signatures}

For the Euclidean spectral analysis use
\begin{equation}
X_\alpha=C_\alpha\times\RR^2,\qquad
\dd s_E^2=\dd r^2+\alpha^2r^2\dd\varphi^2
 +\dd y_0^2+\dd y_3^2,
\quad
0\le\varphi<2\pi,
\quad 0<\alpha\le1.
\label{eq:geometry}
\end{equation}
The singular stratum is the codimension-two set
\[
\WW=\{r=0\}\times\RR^2.
\]
Its transverse measure is
\[
\dd V_\perp=\alpha r\,\dd r\,\dd\varphi.
\]

For the Lorentzian polynomial catalogue the tangent signature is
\(\eta=\diag(1,-1,-1,-1)\), and
\begin{equation}
P^2=q_0^2-q_3^2.
\label{eq:P2}
\end{equation}
The symbols \(q_0,q_3\) first denote the two polynomial variables
printed in v4. Their conversion to covariant or contravariant
momentum components must be fixed when a Ward contraction is made;
Section~\ref{sec:ghost} does so explicitly.

The apex is a prescribed conical worldsheet, not an event horizon.
No state-dependent metric update is present
\cite{CheungIB,CheungIBG}.

\subsection{Scalar and helicity realizations}

A scalar channel with fractional twist \(b\in[0,1)\) has radial
Bessel orders
\[
\nu_n=\frac{\abs{n+b}}{\alpha},\qquad n\in\ZZ.
\]
The Friedrichs realization chooses the regular branch
\(r^{\nu_n}\), with the no-logarithm branch at \(\nu_n=0\).

For a geometric helicity \(S\), the fractional label is
\begin{equation}
b_\Fr(S)=\{S\alpha\},
\label{eq:fractionallabel}
\end{equation}
with fractional part in \([0,1)\).
In the interval \(1/2<\alpha<1\), the distinct labels for
\(S=0,\pm1,\pm2\) are
\[
0,\quad \alpha,\quad 1-\alpha,\quad
2\alpha-1,\quad 2-2\alpha.
\]

For the CR comparison, the v4 labels are
\begin{equation}
b_\CR(S)=
\begin{cases}
S\alpha,&S\ge0,\\
1+S\alpha,&S<0.
\end{cases}
\label{eq:CRlabels}
\end{equation}
For \(S=\pm2\) and \(1/2<\alpha<1\), these labels lie outside
the fractional interval. They are not ordinary Aharonov--Bohm
holonomies modulo one. They encode an additional endpoint choice.

Put
\begin{equation}
\nu=2-\alpha^{-1}\in(0,1).
\label{eq:nu}
\end{equation}
In each exceptional weight-2 channel, CR replaces the regular
Bessel realization of order \(\nu\) by the scale-invariant singular
realization with generalized eigenfunction \(J_{-\nu}\).
Both are nonnegative self-adjoint Bessel realizations in this range
\cite{DerezinskiRichard}.
Their difference, not a periodic relabeling alone, produces the
heat shift proved in Section~\ref{sec:flip}.

\begin{remark}[Parameter range of the printed polynomials]
\label{rem:range}
The common rational forms used below describe the sector
\(1/2<\alpha<1\), with their continuous values at \(\alpha=1/2\).
The F expressions also have their ordinary flat value at \(\alpha=1\).

For \(0<\alpha<1/2\), the appropriate Friedrichs weight-2 label is
\(2\alpha\), not \(2\alpha-1\).
Thus the displayed F polynomials must not be advertised as one
unchanged formula on every \(0<\alpha\le1\).
The scalar heat theorem itself holds for every positive \(\alpha\).
\end{remark}

\subsection{CR selection is not a physical core limit}

The supplied I-B-G-S excerpts describe CR as a mathematical domain
selection restoring the stated source-test de Donder compatibility
and packet curvature regularity in the exceptional sector
\cite[{\loc{def:geo}, \loc{thm:selection},
\loc{thm:spin2}}]{CheungIBGS}.
That selection is not a physical smooth-core theorem.

The supplied I-B-G-S-B description reports that its declared
nonnegative tensor connection-Laplacian and Lichnerowicz-form
families, with the metric spin connection, select a different
positive-core limit. In its neutral exceptional spin-2 channels
for \(\alpha>1/2\), the description excludes the idealized CR
branch in the Mosco/strong-resolvent setting and retains the
corresponding source-test defect
\cite[{\loc{cor:exclude}, \loc{thm:contrast},
\loc{prop:warddefect}}]{CheungTensor}.
This is a source-qualified comparison; a complete independent
tensor proof is not asserted here.

Consequently, CR remains an input to this computation.
Analytical angular-tail, complex-kernel, and Cauchy estimates
are also distinct from an executed, certified full-kernel run
\cite[{\loc{lem:modetail}, \loc{lem:complexbound},
\loc{thm:Cauchy}}]{CheungTensor}.
\separation

\subsection{Constant transverse profile is a jet restriction}

A constant-profile external field is understood as a local jet:
it is constant in the transverse variables near the apex, with
longitudinal Fourier dependence. It need not be a globally
square-integrable constant section of an infinite cone.

Only helicity-zero constant profiles are invariant under the
transverse identification. For a symmetric tensor, the permitted
transverse scalar is the trace combination
\[
h_{11}=h_{22}=f_0.
\]
The individual coordinates \(h_{11}\) and \(h_{22}\) are not
separately helicity-zero eigencomponents.

For comparison with v4, the component keys are
\[
0\leftrightarrow00,\quad
3\leftrightarrow03,\quad
9\leftrightarrow33,\quad
4\leftrightarrow11,\quad
7\leftrightarrow22.
\]
The supplied component array may retain entries in these coordinates;
a genuine constant-profile scalar test uses the transverse trace
combination. In particular, if \(B_{44}=B_{77}\),
\begin{equation}
B_{f_0f_0}=B_{44}+2B_{47}+B_{77}
          =2(B_{44}+B_{47}).
\label{eq:tracecombination}
\end{equation}
No omitted non-lattice component is inferred to vanish.

%======================================================================
\section{A normalized apex heat coefficient}
\label{sec:heat}
%======================================================================

\subsection{Why an ordinary trace is unavailable}

The heat operator on the infinite cone has infinite ordinary trace.
Its bulk volume term must be subtracted before an apex coefficient
is stated. We use an explicit radial Abel subtraction.

Let \(A_b^\Fr\) be the nonnegative twisted Friedrichs transverse
Laplacian. Its diagonal heat kernel is
\begin{equation}
K_b^\Fr(s;r,\varphi;r,\varphi)
 =
\frac{e^{-z}}{4\pi\alpha s}
 \sum_{n\in\ZZ} I_{\abs{n+b}/\alpha}(z),
\qquad
z=\frac{r^2}{2s},
\quad s>0.
\label{eq:heatkernel}
\end{equation}
This is the Bessel heat representation, with normalization relative
to \(\dd V_\perp\).
The required Bessel integral formulas are standard
\cite{DLMF,DerezinskiRichard}; the subtraction and constant are
computed below.

\begin{definition}[Abel apex trace]
\label{def:apex}
For \(u>0\), define
\begin{align}
T_b(u)
 &=\int_{C_\alpha}
 e^{-(\cosh u-1)r^2/(2s)}
 K_b^\Fr(s;x,x)\,\dd V_\perp,\nonumber\\
T_{\rm bulk}(u)
 &=\int_{C_\alpha}
 e^{-(\cosh u-1)r^2/(2s)}
 \frac{\dd V_\perp}{4\pi s}
 =\frac{\alpha}{2(\cosh u-1)}.
\label{eq:Abeltrace}
\end{align}
When the limit exists, put
\[
\ApTr e^{-sA_b^\Fr}
 =\lim_{u\downarrow0}\bigl(T_b(u)-T_{\rm bulk}(u)\bigr).
\]
This definition specifies a bulk-subtracted scale-invariant
coefficient, not an ordinary trace on \(L^2(C_\alpha)\).
\end{definition}

\begin{theorem}[Twisted Friedrichs apex constant]
\label{thm:apex}
For every \(\alpha>0\), \(0\le b<1\), and \(s>0\),
\begin{equation}
\ApTr e^{-sA_b^\Fr}
 =\Ccb_b^\Fr
 :=\frac1{12}\left(\frac1\alpha-\alpha\right)
   -\frac{b(1-b)}{2\alpha}.
\label{eq:Cb}
\end{equation}
The result is independent of the proper time \(s\).
For
\[
f(\lambda)=\int_a^b w(s)e^{-s\lambda}\,\dd s,
\qquad 0<a<b<\infty,\quad w\in L^1([a,b]),
\]
the corresponding heat-superposition apex functional satisfies
\begin{equation}
\ApTr f(A_b^\Fr)=\Ccb_b^\Fr f(0).
\label{eq:functional}
\end{equation}
The subscript \(b\) on the operator denotes the twist;
the integration endpoints in this statement are merely positive
proper-time bounds.
\end{theorem}

\begin{proof}
Integrating \eqref{eq:heatkernel} in the angular and radial variables
gives
\[
T_b(u)=\frac12\sum_{n\in\ZZ}
 \int_0^\infty e^{-z\cosh u}
 I_{\abs{n+b}/\alpha}(z)\,\dd z.
\]
For \(u>0\) and \(\rho>-1\),
\[
\int_0^\infty e^{-z\cosh u}I_\rho(z)\,\dd z
 =\frac{e^{-\rho u}}{\sinh u}.
\]
The nonnegative summands justify the interchange.
The geometric series gives
\[
\sum_{n\in\ZZ}e^{-\abs{n+b}u/\alpha}
 =
\frac{\cosh((1/2-b)u/\alpha)}{\sinh(u/(2\alpha))}.
\]
Hence
\begin{equation}
T_b(u)=
\frac{\cosh((1/2-b)u/\alpha)}
 {2\sinh u\,\sinh(u/(2\alpha))}.
\label{eq:closedAbel}
\end{equation}

With \(B_2(b)=b^2-b+1/6\), expansion at zero yields
\[
T_b(u)
 =\frac{\alpha}{u^2}
  +\frac{B_2(b)}{2\alpha}-\frac{\alpha}{6}
  +O(u^2),
\]
whereas
\[
T_{\rm bulk}(u)
 =\frac{\alpha}{u^2}-\frac{\alpha}{12}+O(u^2).
\]
Subtracting proves
\[
\Ccb_b^\Fr
 =\frac{B_2(b)}{2\alpha}-\frac{\alpha}{12},
\]
which is \eqref{eq:Cb}.

The radial change of variable removed \(s\), proving exact
proper-time independence. For a compact proper-time interval,
the closed expressions and their subtracted limits are uniformly
bounded. Dominated convergence and linearity then give
\[
\ApTr f(A_b^\Fr)
 =\Ccb_b^\Fr\int_a^b w(s)\,\dd s
 =\Ccb_b^\Fr f(0).
\]
\end{proof}

The untwisted coefficient and conical heat-trace methodology are
classical \cite{Fursaev1994,FursaevMiele1997,Vassilevich}.
The proof fixes the normalization rather than attributing a
numerical re-derivation to unavailable software.

\begin{remark}[Function class]
Theorem~\ref{thm:apex} does not assert an unrestricted ordinary-trace
identity for every Schwartz function.
Extending \eqref{eq:functional} to resolvent products,
differentiated kernels, or meromorphic graph families requires
the corresponding functional-calculus and interchange hypotheses.
Those hypotheses are stated in Section~\ref{sec:reduction}.
\end{remark}

\subsection{Apex coefficients are not only diffracted-wave coefficients}

The scalar Sommerfeld--Carslaw representation, extension-dependent
Green identities, and supported transfer in I-B-G concern continuum
operators on their stated domains
\cite[{\loc{thm:green}, \loc{eq:Carslaw},
\loc{thm:transfer}}]{CheungIBG}.
They do not identify a loop residue with a finite detector output.

One further distinction is necessary: the apex heat constant can
receive contributions from geometric images even when a chosen
diffracted wave term vanishes.

\begin{example}[A nonzero apex constant in a pure-image case]
\label{ex:image}
For the neutral cone \(\alpha=1/2\),
\[
\Ccb_0^\Fr=\frac18.
\]
The cone is a two-fold rotation quotient. The nontrivial image has
plane heat kernel
\[
\frac1{4\pi s}\exp(-r^2/s)
\]
on the diagonal after a half-turn.
Integrating over the wedge of angle \(\pi\) gives
\[
\frac{\pi}{4\pi s}
 \int_0^\infty r e^{-r^2/s}\,\dd r=\frac18.
\]
Thus this apex coefficient is present without needing a separately
nonzero diffracted-front term.
\end{example}

The finite-window echo ratios of I-B-ECHO and I-C concern a
different functional. Their additive front expansion and their
zero-deficit witness on \(1/2<\alpha\le1\) do not turn
\(\Ccb_b\) into that ratio
\cite{CheungEcho,CheungIC}.
\separation

%======================================================================
\section{Singular Bessel channels and the CR heat shift}
\label{sec:flip}
%======================================================================

For \(0<\nu<1\), let \(A_\nu^+\) and \(A_\nu^-\) denote the
scale-invariant regular and singular nonnegative self-adjoint radial
realizations in \(L^2((0,\infty),r\,\dd r)\).
Their radial heat kernels are
\[
k_\nu^\pm(s;r,r')
 =\frac1{2s}
  e^{-(r^2+r'^2)/(4s)}
  I_{\pm\nu}\!\left(\frac{rr'}{2s}\right).
\]
These choices are endpoint realizations, not ordinary potentials
added to the regular cone.

\begin{theorem}[One-channel flip]
\label{thm:flip}
For \(0<\nu<1\),
\begin{equation}
\Tr_{\rm rad}
 \bigl(e^{-sA_\nu^-}-e^{-sA_\nu^+}\bigr)=\nu,
\qquad s>0.
\label{eq:flip}
\end{equation}
The diagonal relative integral is convergent.
Consequently, replacing one such regular angular channel by its
singular realization increases its apex heat constant by \(\nu\).
\end{theorem}

\begin{proof}
The radial trace difference is
\[
\frac12\int_0^\infty
 e^{-z}\bigl(I_{-\nu}(z)-I_\nu(z)\bigr)\,\dd z.
\]
Near zero, the worst term is \(z^{-\nu}\), which is integrable
because \(\nu<1\).
At infinity use
\[
I_{-\nu}(z)-I_\nu(z)
 =\frac{2\sin(\pi\nu)}{\pi}K_\nu(z).
\]
The resulting integrand decays exponentially
\cite{DLMF}.

Insert the Abel factor and use the Laplace transform from the
preceding proof. The integral becomes
\[
\frac{e^{\nu u}-e^{-\nu u}}{2\sinh u}
 =\frac{\sinh(\nu u)}{\sinh u}.
\]
Dominated convergence gives the limit \(\nu\).
Scale invariance makes it independent of \(s\).
\end{proof}

\begin{corollary}[Three apex coefficients in the comparison sector]
\label{cor:threeC}
For \(1/2<\alpha<1\), define
\begin{align}
c_0&=\frac{1-\alpha^2}{12\alpha},\nonumber\\
c_1&=\frac{5\alpha^2-6\alpha+1}{12\alpha},\nonumber\\
c_2^\Fr&=\frac{23\alpha^2-36\alpha+13}{12\alpha},\nonumber\\
c_2^\CR&=\frac{23\alpha^2-12\alpha+1}{12\alpha}.
\label{eq:cvalues}
\end{align}
These are the per-helicity apex coefficients for weight zero,
weight one, and the two respective weight-two realizations.
They satisfy
\begin{align}
c_1-c_0&=\frac{\alpha-1}{2},\label{eq:c1diff}\\
c_2^\CR-c_0&=2\alpha-1,\label{eq:c2CRdiff}\\
c_2^\Fr-c_0&=
 \frac{(\alpha-1)(2\alpha-1)}{\alpha},\label{eq:c2Fdiff}\\
c_2^\CR-c_2^\Fr&=
 2-\frac1\alpha=\nu.
\label{eq:branchdiff}
\end{align}
At \(\alpha=1/2\), continuous evaluation gives
\[
c_0=c_2^\Fr=c_2^\CR=\frac18,\qquad c_1=-\frac18.
\]
\end{corollary}

\begin{proof}
Apply \eqref{eq:Cb} at \(b=0,\alpha,2\alpha-1\).
For the exceptional CR channel add \(\nu\) from
Theorem~\ref{thm:flip}. Expansion gives \eqref{eq:cvalues};
subtraction gives the remaining identities.
The mirror labels have the same coefficients because
\(b(1-b)=(1-b)b\).
\end{proof}

The polynomial expression obtained by inserting \(b=2\alpha\)
into \eqref{eq:Cb} reproduces \(c_2^\CR\).
This is now justified by a specified channel flip.
It is not a general license to use arbitrary real \(b\) as a
fractional twist.

\begin{remark}[The CR flat-limit remnant]
\label{rem:flatCR}
As \(\alpha\uparrow1\),
\[
c_0\to0,\qquad c_1\to0,\qquad
c_2^\Fr\to0,\qquad c_2^\CR\to1.
\]
The last statement is a one-sided limit of the singular-channel
family. At \(\alpha=1\), its formal order is \(\nu=1\), and
\[
\int_0^\epsilon r\,\abs{r^{-1}}^2\,\dd r=\infty.
\]
The deficient singular \(L^2\) branch used above no longer exists
in the same form. Therefore a nonzero limiting CR polynomial is
not, without an additional theorem, the heat trace of a CR
realization at the exactly smooth endpoint.
Nonuniform spin-dependent flat limits are also discussed in the
established conical heat-kernel literature
\cite{FursaevMiele1997}.
\end{remark}

%======================================================================
\section{An \texorpdfstring{\(r^2\)}{r\textasciicircum 2}-weighted apex coefficient (new in v6)}
\label{sec:weighted}
%======================================================================

Constant transverse profiles see only \(\Ccb_b\).
A quadratic transverse jet \(G=r^2g\) needs the next moment of the
apex heat distribution.
v5 recorded that no value of the higher coefficient \(\mathcal D_b\)
had been supplied or derived.
This section derives it with the same Abel subtraction used in
Definition~\ref{def:apex}.

\begin{definition}[Weighted Abel apex trace]
\label{def:weighted}
For \(u>0\) put
\begin{align*}
T^{(2)}_b(u)
 &=\int_{C_\alpha}r^2
 e^{-(\cosh u-1)r^2/(2s)}K_b^\Fr(s;x,x)\,\dd V_\perp,\\
T^{(2)}_{\rm bulk}(u)
 &=\int_{C_\alpha}r^2
 e^{-(\cosh u-1)r^2/(2s)}\frac{\dd V_\perp}{4\pi s}
 =\frac{\alpha s}{(\cosh u-1)^2}.
\end{align*}
When the limit exists, set
\(\ApTr[r^2e^{-sA_b}]=\lim_{u\downarrow0}
(T^{(2)}_b(u)-T^{(2)}_{\rm bulk}(u))\).
\end{definition}

\begin{theorem}[Weighted twisted Friedrichs apex coefficient]
\label{thm:weighted}
Let \(\alpha>0\), \(0\le b<1\), and \(s>0\).
With Bernoulli polynomials \(B_2(b)=b^2-b+\tfrac16\) and
\(B_4(b)=b^4-2b^3+b^2-\tfrac1{30}\),
\begin{equation}
\ApTr\bigl[r^2e^{-sA_b^\Fr}\bigr]=\mathcal D_b\,s,
\qquad
\mathcal D_b=-\frac{11\alpha}{180}
+\frac{B_2(b)}{3\alpha}-\frac{B_4(b)}{6\alpha^3}.
\label{eq:Db}
\end{equation}
Equivalently,
\begin{equation}
\mathcal D_b=
\frac{1+10\alpha^2-11\alpha^4-60\alpha^2b(1-b)-30b^2(1-b)^2}
{180\alpha^3}.
\label{eq:Dbexpanded}
\end{equation}
For a heat superposition \(f\) as in Theorem~\ref{thm:apex},
\begin{equation}
\ApTr\bigl[r^2f(A_b^\Fr)\bigr]=-\mathcal D_b\,f'(0).
\label{eq:Dfunctional}
\end{equation}
\end{theorem}

\begin{proof}
With \(z=r^2/(2s)\), we have \(r^2=2sz\) and
\(r\,\dd r=s\,\dd z\).
The angular integral and the heat representation
\eqref{eq:heatkernel} then give
\[
T^{(2)}_b(u)=s\sum_{n\in\ZZ}\int_0^\infty
z\,e^{-z\cosh u}I_{\abs{n+b}/\alpha}(z)\,\dd z.
\]
Differentiate the Laplace transform from the proof of
Theorem~\ref{thm:apex} with respect to \(u\):
\[
\int_0^\infty z\,e^{-z\cosh u}I_\rho(z)\,\dd z
=-\frac1{\sinh u}\frac{\dd}{\dd u}
\frac{e^{-\rho u}}{\sinh u},
\qquad \rho>-1.
\]
Since all summands are nonnegative, the sum may be interchanged
with the integral, giving
\[
\frac{T^{(2)}_b(u)}{s}=-\frac1{\sinh u}\frac{\dd}{\dd u}
\left[\frac{\cosh((\tfrac12-b)u/\alpha)}
{\sinh u\,\sinh(u/(2\alpha))}\right].
\]
Laurent expansion at \(u=0\) shows that the
\(u^{-4},u^{-3},u^{-2},u^{-1}\) terms of
\(T^{(2)}_b/s\) and \(T^{(2)}_{\rm bulk}/s=\alpha/(\cosh u-1)^2\)
coincide.
The difference of their constant terms is \eqref{eq:Db}.
Expanding the Bernoulli polynomials gives \eqref{eq:Dbexpanded}.

For a heat superposition \(f(\lambda)=\int w(s)e^{-s\lambda}\dd s\),
dominated convergence on the compact proper-time support gives
\(\ApTr[r^2f(A)]=\mathcal D_b\int s\,w(s)\,\dd s
=-\mathcal D_bf'(0)\).
\textsf{[exec]} The Laurent algebra is carried out in
\loc{code/weighted_apex.py}: all four pole terms vanish identically,
and the constant term equals \eqref{eq:Db}
(\loc{logs/v6/weighted_checks.json}, key \texttt{D\_matches\_log}).
\end{proof}

\begin{proposition}[Weighted singular-channel shift]
\label{prop:weightedflip}
For \(0<\nu<1\), replacing one regular channel of order \(\nu\) by
its singular realization changes the weighted apex coefficient by
\[
\Delta\mathcal D=\frac{2\nu(1-\nu^2)}{3}.
\]
For the CR weight-two channel, \(\nu=2-\alpha^{-1}\) and
\[
\mathcal D^\CR_{S=2}=\mathcal D_{2\alpha-1}+\Delta\mathcal D
=\mathcal D_{b=2\alpha},
\]
the polynomial \eqref{eq:Dbexpanded} evaluated at \(b=2\alpha\).
\end{proposition}

\begin{proof}
The weighted analogue of the proof of Theorem~\ref{thm:flip}
replaces \(e^{-\nu u}\) by \(e^{+\nu u}\) in one channel.
The difference is
\(-(\sinh u)^{-1}\partial_u[2\sinh(\nu u)/\sinh u]\).
Its limit as \(u\to0\) is
\(-u^{-1}\partial_u[2\nu(1+(\nu^2-1)u^2/6)]
\to2\nu(1-\nu^2)/3\).
Near zero the singular kernel is \(O(z^{1-\nu})\), so the weighted
diagonal difference is absolutely integrable.
\textsf{[exec]} The identity
\(\mathcal D_{2\alpha}-\mathcal D_{2\alpha-1}=\Delta\mathcal D\) is
checked exactly as a rational function of \(\alpha\)
(\loc{logs/v6/weighted_checks.json}).
\end{proof}

\begin{corollary}[Comparison-sector values]
\label{cor:Dvalues}
For \(1/2<\alpha<1\):
\begin{align*}
\mathcal D_0&=\frac{(1-\alpha^2)(1+11\alpha^2)}{180\alpha^3},\\
\mathcal D_{S=1}&=\frac{(\alpha-1)(\alpha+1)(19\alpha^2-1)}{180\alpha^3},\\
\mathcal D^\Fr_{S=2}&=-\frac{(\alpha-1)(251\alpha^3-829\alpha^2+601\alpha-119)}{180\alpha^3},\\
\mathcal D^\CR_{S=2}&=-\frac{(\alpha-1)(251\alpha^3-109\alpha^2+\alpha+1)}{180\alpha^3}.
\end{align*}
All four vanish as \(\alpha\uparrow1\).
In particular, the weighted CR coefficient has \emph{no} flat
remnant, unlike \(c_2^\CR\to1\) (Remark~\ref{rem:flatCR}),
because \(\Delta\mathcal D\to0\) as \(\nu\to1\).
\end{corollary}

\begin{proposition}[Orbifold checks]
\label{prop:orbifold}
For \(\alpha=1/N\) with \(N\in\{2,\dots,8\}\),
\[
\mathcal D_0(1/N)=\frac1{4N}\sum_{k=1}^{N-1}\csc^4\frac{\pi k}{N}
=\frac{(N^2-1)(N^2+11)}{180N}.
\]
This gives \(1/8,\,8/27,\,9/16,\,24/25,\,329/216,\,16/7,\,105/32\).
For the twisted orbifolds \((N,j)=(3,1),(4,1),(5,2)\), with
\(b=j/N\), the image sums weighted by holonomy phases equal
\(\mathcal D_b\), namely \(-4/27\), \(-1/16\), and \(-16/25\).
\end{proposition}

\begin{proof}
On \(\RR^2/\ZZ_N\) the heat kernel is the image sum.
The \(k\)-th image contributes
\((4\pi s)^{-1}e^{-r^2\sin^2(\pi k/N)/s}\).
Integrating \(r^2\) times it over the wedge of angle \(2\pi/N\) gives
\(s/(4N\sin^4(\pi k/N))\).
The identity element is the bulk term.
\textsf{[exec]} Both sides are evaluated exactly for each \(N\),
and to \(10^{-12}\) for the twisted cases.
See the key \texttt{orbifold\_csc4} in
\loc{logs/v6/weighted_checks.json}, and
\loc{logs/v6/weighted_apex.json}.
\end{proof}

\begin{remark}[Literature comparison]
The untwisted moment \(\mathcal D_0\) has the
\(\csc^4\)-sum structure of the classical conical heat coefficients
\cite{Fursaev1994}.
The identification of \(2\pi\alpha\,\mathcal D_0\) with the
\(C_4\)-type coefficient of that reference is made from memory of its
normalization (\(\dagger\)). It is not used as an input: the
orbifold identity above is self-contained.

Two literature results are analogues of the CR one-sided remnant
of \(c_2^\CR\): the non-uniform spin-dependent flat limits of
\cite{FursaevMiele1997}, and the spin-one horizon contact term of
\cite{Kabat}.
Corollary~\ref{cor:Dvalues} shows that this remnant does
\emph{not} propagate to the \(r^2\) moment.
No numerical spin-2 coefficient is imported from either work.
\end{remark}

\subsection{Application: the second \texorpdfstring{\(\phi\phi\)}{phi-phi} combination (partial)}

By \eqref{eq:Dfunctional}, a quadratic transverse jet contributes
\(-\mathcal D_S\) times the derivative, at zero transverse spectral
parameter, of the longitudinal integrand.
Within Assumption~\ref{ass:reduction}:
\begin{itemize}
\item \textsf{[exec]} \emph{D05a.} After the weighted insertion,
the single-line tadpole has longitudinal integrand proportional to
\(p^2\mathcal D_S/l^4\).
At \(D=2\) this is \(O(R^{-4})\) at large radius, so it has no
logarithmic UV term (Lemma~\ref{lem:UV}).
The D05a contribution to the quadratic-profile pole is therefore
\emph{computed} to be \(0\).
\item \emph{D04.} The bubble couples a graviton line and a scalar
line in different channels.
Its \(r^2\) insertion requires mixed-channel traces of the
transverse position operators \(x_\pm\) between channels \(S\) and
\(S\pm1\).
These are not functions of a single \(A_S\), so they fall outside
Theorem~\ref{thm:weighted} and outside (R2).
The D04 part is \(\OPEN\).
\end{itemize}
The second scalar basis value is therefore \emph{partial}.
It is not reported as equal to its D05a part.

%======================================================================
\section{Conditional spectral reduction of a two-point graph}
\label{sec:reduction}
%======================================================================

\subsection{Helicity classes and one closed transverse trace}

Under longitudinal/transverse splitting, a symmetric four-dimensional
tensor has:
\begin{itemize}
\item four weight-zero components: the symmetric longitudinal block
and the transverse trace;
\item four weight-one components: the longitudinal--transverse block;
\item two weight-two components: the transverse traceless block.
\end{itemize}
A vector has two weight-zero and two weight-one components.
A scalar has weight zero.

These are geometric helicities. The older labels \(L,M,T\) do not,
without their definitions, determine the correct apex weight.

A reduced closed loop has one transverse spectral trace.
Accordingly, a contribution normalized to one common spectral
channel receives one coefficient \(c_S\).
It does not receive the product of one apex coefficient for every
propagator in that trace.

\subsection{The missing hypotheses}

\begin{assumption}[Graph-specific spectral reduction contract]
\label{ass:reduction}
For a specified graph \(G\), external constant helicity-zero jet,
and fixed routing, require:
\begin{enumerate}[label=\textup{(R\arabic*)}]
\item \textbf{Complete vertex data.}
The tensor vertices, propagator numerators, source insertions,
statistics signs, symmetry factors, and independent-component
normalizations are specified.

\item \textbf{Operator reduction.}
After the complete vertex contraction, each contributing transverse
operator expression is reduced to a finite sum
\[
f_{G,S}(A_S;l,p),
\]
on one identified channel domain, plus a remainder proved not
to contribute to the apex UV residue.

\item \textbf{Domain and trace legitimacy.}
Every differentiated intertwiner, integration by parts, and trace
rearrangement is valid on that domain.
Any apex boundary term or finite-part trace defect is included,
or proved to vanish after the full permitted sum.

\item \textbf{UV/IR separation.}
An auxiliary infrared prescription, such as a positive mass in
the Schwinger representation, distinguishes the UV residue
from massless infrared poles.
A scaleless integral is not assigned its UV pole merely by
ignoring its infrared part.

\item \textbf{Meromorphic interchange.}
The heat-superposition apex identity extends to these regulated
spectral functions. The transverse subtraction, longitudinal
integration, removal of auxiliary cutoffs, and Laurent-residue
operation commute in the stated meromorphic prescription.
\end{enumerate}
A graph is called \emph{spectrally certified in this sector}
only after these conditions have been established for that graph.
\end{assumption}

\begin{theorem}[Conditional spectral apex reduction]
\label{thm:reduction}
Under Assumption~\ref{ass:reduction}, the apex UV residue of \(G\)
is
\begin{equation}
\ResUV G_{\rm ap}(p)
 =
\sum_S c_S\,
 \ResUV\!
 \left[
 \mu^{4-d}
 \int\frac{\dd^{\,d-2}l_\parallel}{(2\pi)^{d-2}}
 f_{G,S}(0;l_\parallel,p)
 \right].
\label{eq:reduction}
\end{equation}
Multiplicity, statistics, coupling, and vertex phases are included
in the corresponding \(f_{G,S}\).
For a flat-integrand representation this is the
\(l_\perp=0\) expression only when the operator reduction
\textup{(R2)--(R3)} has actually established that identification.
\end{theorem}

\begin{proof}
First retain the auxiliary proper-time and infrared cutoffs.
By (R2), the complete graph is a finite sum of functions of the
identified \(A_S\), up to the residue-free remainder.
Apply \eqref{eq:functional} and Theorem~\ref{thm:flip} to each
regularized heat-superposition term. This replaces its transverse
apex trace by \(c_S f_{G,S}(0;l,p)\).

Condition (R3) ensures that the transformations used to reach
these terms have not omitted a boundary or trace defect.
Condition (R5) then allows longitudinal integration, cutoff removal,
and residue extraction in that order or the stated equivalent
order. The remainder contributes zero by (R2).
This gives \eqref{eq:reduction}.
\end{proof}

\begin{remark}[What the heat theorem does not prove]
A bubble containing two different channel realizations, a
domain-changing differential operator, or an unaccounted boundary
term is not automatically a function of one \(A_S\).
The phrase ``the external transverse momentum is zero'' does not
establish (R2) or (R3).
Theorem~\ref{thm:reduction} therefore cannot be replaced by the
unqualified statement that every differentiated two-point graph
equals its flat integrand at \(l_\perp=0\).
\end{remark}

\subsection{Intertwining is a domain statement}

The formal Bessel recurrences give differential identities between
solutions. Equality of self-adjoint resolvents additionally requires
the differentiated solution to lie in the target domain
\cite{DLMF,DerezinskiRichard}.

For example, if a singular leading term behaves as \(r^{-\nu}\),
one transverse derivative can produce \(r^{-\nu-1}\).
For \(\nu>0\),
\[
\int_0^\epsilon
 r\,\abs{r^{-\nu-1}}^2\,\dd r
 =\int_0^\epsilon r^{-2\nu-1}\,\dd r=\infty.
\]
Thus a formal recurrence is not an unrestricted \(L^2\) intertwiner
through the next helicity.

The supplied v4 record reports:
\begin{itemize}
\item D03, D04, and D11 as clean in its constant-profile reduction;
\item D06 and D13 as clean on CR;
\item D01 as conditional on \(\HI\), because its derivative
tower can reach the \(S=2\to3\) sector.
\end{itemize}
These are retained as graph-assessment reports, not replaced by
an unavailable local derivation.

For D01, the unrestricted tower version of \(\HI\) is not an
available hypothesis: the same record states its failure.
A nonvacuous successor must prove that the \emph{actual contracted
graph} avoids the forbidden terms, includes their boundary
contributions, or cancels them in the complete sum.
We denote that graph-local requirement by
\[
(\mathrm H_{I,\mathrm{D01}}).
\]
Merely assuming the already excluded unrestricted tower would not
certify D01.

\subsection{The D01 obstruction made exact (new in v6)}
\label{sec:HI}

Let \(\nu=2-\alpha^{-1}\in(0,1)\).
In the flipped weight-two channel, the CR--F difference of radial
heat kernels is
\[
\Delta k(s;r,r')=k^-_\nu-k^+_\nu
=\frac{1}{2s}e^{-(r^2+r'^2)/(4s)}\,
\frac{2\sin\pi\nu}{\pi}\,K_\nu\!\left(\frac{rr'}{2s}\right).
\]
The D01 derivative tower carries this channel into \(S=3\) through
the raising operator
\[
X=\partial_r-\frac{\nu}{r},
\]
the radial part of the helicity-raising derivative on the
\(J_{-\nu}\) branch.
The term that \((\mathrm H_{I,\mathrm{D01}})\) must control is the
composite diagonal
\[
\mathcal R(s;r)=\bigl[X_rX_{r'}\Delta k(s;r,r')\bigr]_{r'=r}.
\]

\begin{theorem}[Exact obstruction and Mellin value for D01]
\label{thm:HI}
\leavevmode
\begin{enumerate}[label=\textup{(\roman*)}]
\item \emph{Obstruction.}
As \(r\to0\), \(\mathcal R(s;r)=c_*\,r^{-2\nu-2}+O(r^{-2\nu})\),
with \(c_*\neq0\) for every \(0<\nu<1\).
Hence \(\int_0^\epsilon r\,\mathcal R\,\dd r\) diverges.
No \(L^2\) intertwiner and no trace-class statement holds for the
composite. The unrestricted \(\HI\) fails, exactly as the v4 record
states.

\item \emph{Mellin value.}
The Mellin transform
\(M(z)=\int_0^\infty r^{1+z}\mathcal R(s;r)\,\dd r\)
converges for \(\Re z>2\nu\).
It continues meromorphically, with poles contained in
\((2\nu-2\NN_0)\cup(-2\nu-2\NN_0)\).
So it is regular at \(z=0\) for \(0<\nu<1\), and
\[
M(0)=0\qquad\text{identically in }\nu\in(0,1).
\]
\end{enumerate}
\end{theorem}

\begin{proof}
Put \(s=1/2\), so that \(x=r^2\); by scale invariance
\(M(0)\) scales as \(s^{-1}\).
Then
\[
\mathcal R=A\,e^{-x}\Bigl[xK-2xK'+K'+xK''
-2\nu(K'-K)+\nu^2x^{-1}K\Bigr],
\qquad K=K_\nu(x),
\]
where \(A=\sin(\pi\nu)/(\pi s)\).
Use \(K'=-\tfrac12(K_{\nu-1}+K_{\nu+1})\) and
\(K''=\tfrac14(K_{\nu-2}+2K_\nu+K_{\nu+2})\).
Each term then has the form
\(\int_0^\infty x^{\mu-1}e^{-x}K_m(x)\,\dd x\) with
\[
\int_0^\infty x^{\mu-1}e^{-x}K_m(x)\,\dd x
=\frac{\sqrt\pi\,\Gamma(\mu+m)\Gamma(\mu-m)}
{2^\mu\,\Gamma(\mu+\tfrac12)},
\qquad \Re\mu>\abs{\Re m}.
\]
This is the \(a\to1\) limit of DLMF~10.43.22 \cite{DLMF}.
\textsf{[exec]} It is checked by quadrature in its strip to
\(10^{-16}\) (\loc{logs/v6/hi_d01_exact.json}).
The right side is meromorphic in \(\mu\), which gives the stated
continuation and poles.
Leading-order expansion of \(K_\nu\) gives (i).

At \(z=0\) the bracket is a sum of ten gamma ratios with
\(\mu\in\{0,1,2\}\) and \(m\in\{\nu,\nu\pm1,\nu\pm2\}\).
\textsf{[exec]} Applying \(\Gamma(w+1)=w\Gamma(w)\) reduces it
to \(0\) as a symbolic function of \(\nu\)
(\loc{code/hi_d01_exact.py}).
The sum at \(z=10^{-30}\) is \(O(10^{-30})\) for seven values of
\(\nu\).

An independent route uses direct quadrature with the
\(r^{-2\nu-2}\) and \(r^{-2\nu}\) heads subtracted and continued
analytically (\loc{code/hi_d01.py}, 60 digits).
It gives \(s\,M(0)\in[-1.8\times10^{-21},\,3.7\times10^{-32}]\)
at \(\alpha=0.6,0.7,0.75,0.8,0.9\), and the value is independent
of \(s\).
The residual grows like \(\epsilon^{4-2\nu}\) for the cutoff
\(\epsilon=10^{-9}\), consistent with zero.
\end{proof}

\begin{corollary}[Status of \((\mathrm H_{I,\mathrm{D01}})\)]
\label{cor:HI}
Interpret the apex coefficient of the forbidden D01 composite by
homogeneous (Mellin) continuation.
This is the unique scale-covariant extension of the
non-integrable diagonal here, since there is no resonance for
\(0<\nu<1\).
With that reading the composite contributes \(0\), and the D01
constant-profile entries of Table~\ref{tab:hhC} need no correction
from it.

The hypothesis is therefore \emph{upgraded} in two respects: its
failure as an \(L^2\) statement is now exact, and the Mellin value
is exactly zero.
One thing is not shown: that the XVI-P prescription requires
homogeneous continuation of this composite, as opposed to some
other extension of the divergent apex integral.
That is a prescription statement, part of (R3)--(R5), and it is
recorded as \CONDITIONAL.
\end{corollary}

\subsection{Longitudinal UV normalization}

Set
\[
\eps=4-d,\qquad D=d-2=2-\eps.
\]
This \(\eps\) is twice the common convention
\(\eps_{\rm DR}=(4-d)/2\).

\begin{lemma}[Angular extraction of the logarithmic UV residue]
\label{lem:UV}
Suppose an infrared-regular Euclidean longitudinal integrand has
large-radius expansion
\[
I(R\omega;p)=R^{-2}g(\theta;p)+O(R^{-2-\delta}),
\qquad \delta>0,
\]
uniformly in angle, with no other UV-logarithmic term.
If its dimensional continuation is regular in angle at \(D=2\),
the coefficient of \(1/\eps\) is
\begin{equation}
\frac1{(2\pi)^2}
 \int_0^{2\pi}g(\theta;p)\,\dd\theta.
\label{eq:UVangle}
\end{equation}
An overall Wick phase \(+\ii\), when prescribed by the amplitude
convention, multiplies this coefficient by \(\ii\).
\end{lemma}

\begin{proof}
Choose a fixed large radius \(R_0\) and separate the bounded
infrared region. The leading UV radial integral is
\[
\int_{R_0}^\infty R^{D-1}R^{-2}\,\dd R
 =\int_{R_0}^\infty R^{-1-\eps}\,\dd R
 =\frac{R_0^{-\eps}}{\eps}.
\]
The remainder is UV integrable near \(\eps=0\).
The measure at \(D=2\) contributes \((2\pi)^{-2}\dd\theta\).
An \(O(\eps)\) correction to the angular continuation affects
the finite part of a simple pole, not this leading coefficient.
\end{proof}

The factor in \eqref{eq:UVangle} agrees with the v4 expression
\(\ii(4\pi^2)^{-1}\int g\) when its \(+\ii\) Wick convention is
applicable. Graph-specific additional phases are not fixed by
this radial argument.

Dimensional reduction requires a consistent distinction between
momentum and vertex-index spaces; it does not automatically establish
a Ward identity for the present boundary problem
\cite{Stockinger}.

%======================================================================
\section{The retained constant-profile pole catalogue}
\label{sec:catalogue}
%======================================================================

\subsection{Status and notation}

Write \(\widehat B_G\) for the candidate coefficient of
\(1/(4-d)\) printed in v4.
The hat distinguishes a supplied graph record from an independently
derived loop amplitude.

All formulas below are exact algebraic expressions.
\textsf{[exec]} In v6 every entry is checked against the executed
spectral logs of the frozen vertex arrays
(\loc{code/catalogue_check.py}, \loc{logs/v6/catalogue_check.json}).
Every graph, branch, and key agrees exactly, the entries absent from
the logs are computed zeros, and the new entries are marked.
Their interpretation as complete graph residues remains subject to
Assumption~\ref{ass:reduction}.

For either branch, put
\begin{equation}
c_2=
\begin{cases}
c_2^\Fr,&\text{F},\\
c_2^\CR,&\text{CR},
\end{cases}
\quad
A=c_0+2c_1+4c_2,
\quad
B=-c_0+2c_1.
\label{eq:AB}
\end{equation}
The \(c\)'s are the per-channel quantities
\eqref{eq:cvalues}; any multiplicity is already included in the
coefficients of the following tables.

\subsection{Scalar two-point sector}

\begin{table}[htbp]
\centering
\small
\begin{tabularx}{\linewidth}{@{}L{.20\linewidth}Y@{}}
\toprule
Graph & Retained candidate residue\\
\midrule
D04 &
\(\displaystyle -\frac{\ka^2}{4\pi}c_0P^2\)\\[6pt]
D05a &
\(\displaystyle -\frac{\ka^2}{4\pi}(c_2-c_0)P^2\)\\[6pt]
D04+D05a &
\(\displaystyle -\frac{\ka^2}{4\pi}c_2P^2\)\\
\bottomrule
\end{tabularx}
\caption{Equivalent coefficient form of the v4
constant-profile \(\phi\phi\) catalogue.}
\label{tab:phiC}
\end{table}

Expansion gives exactly
\begin{align}
\widehat B_{\mathrm{D04}}
 &=\frac{\ka^2(\alpha^2-1)}{48\pi\alpha}P^2,\nonumber\\
\widehat B_{\mathrm{D05a}}^\CR
 &=-\frac{\ka^2(2\alpha-1)}{4\pi}P^2,\nonumber\\
\widehat B_{\mathrm{D05a}}^\Fr
 &=-\frac{\ka^2(\alpha-1)(2\alpha-1)}
 {4\pi\alpha}P^2,\nonumber\\
\widehat B_{\mathrm{D04+D05a}}^\CR
 &=\frac{\ka^2(-23\alpha^2+12\alpha-1)}
 {48\pi\alpha}P^2,\nonumber\\
\widehat B_{\mathrm{D04+D05a}}^\Fr
 &=\frac{\ka^2(\alpha-1)(13-23\alpha)}
 {48\pi\alpha}P^2.
\label{eq:phiExpanded}
\end{align}

The one-sided CR limit of the sum as \(\alpha\uparrow1\) is
\(-\ka^2P^2/(4\pi)\).
It is the polynomial remnant discussed in
Remark~\ref{rem:flatCR}, not a claimed flat-endpoint operator value.

\textsf{[exec]} The D05a symmetry factor \(S_f=1/2\) is read directly
from the regenerated Q02 graph record
\loc{xviq_regen/Q02_bulk_graphs.json} \cite{CheungQ}.
The same record gives \(1/2\) for D12 and \(1\) for D04.
With the XVI-P one-propagator tadpole rule, the engine prefactor is
\(-\ka^2/2\), and the displayed D05a values follow without a fitted
constant.

\subsection{Tensor two-point sector}

For the next table, factor out \(\ka^2/(16\pi)\).
Indices denote the independent-coordinate entries listed in
Section~\ref{sec:conventions}.

\begin{longtable}{@{}L{.27\linewidth}L{.64\linewidth}@{}}
\caption{The \(hh\) bubble entries of v4, in equivalent
apex-coefficient form, with the \((5,5)\) entries added from the
executed record (v6).
Every displayed entry is multiplied by \(\ka^2/(16\pi)\).
D01 retains the graph-local intertwining qualification.}
\label{tab:hhC}\\
\toprule
Graph and key & Coefficient after removing \(\ka^2/(16\pi)\)\\
\midrule
\endfirsthead
\toprule
Graph and key & Coefficient after removing \(\ka^2/(16\pi)\)\\
\midrule
\endhead
\bottomrule
\endfoot

D11:
\((4,4),(4,7),(7,7)\)
&
\(\displaystyle \frac12 c_0P^2\)\\[5pt]

D03:
\((0,0),(0,9),(9,9)\)
&
\(\displaystyle -c_0P^2\)\\[5pt]

D03:
\((3,3)\)
&
\(\displaystyle 4c_0P^2\)\\[5pt]

D03:
\((0,4)\)
&
\(\displaystyle c_0(q_0^2-3q_3^2)\)\\[5pt]

D03:
\((3,4)\)
&
\(\displaystyle 4c_0q_0q_3\)\\[5pt]

D03:
\((4,9)\)
&
\(\displaystyle -c_0(3q_0^2-q_3^2)\)\\[5pt]

D03:
\((4,4)\)
&
\(\displaystyle -(c_0+4c_1)P^2\)\\[5pt]

D03:
\((4,7)\)
&
\(\displaystyle -c_0P^2\)\\[5pt]

D03:
\((5,5)\) \emph{(new in v6)}
&
\(\displaystyle -8c_1P^2\)\\[5pt]

D01:
\((0,0),(0,9),(9,9)\)
&
\(\displaystyle c_0P^2\)\\[5pt]

D01:
\((3,3)\)
&
\(\displaystyle -4c_0P^2\)\\[5pt]

D01:
\((4,4)\)
&
\(\displaystyle (c_0-4c_1+c_2)P^2\)\\[5pt]

D01:
\((4,7)\)
&
\(\displaystyle (c_0+c_2)P^2\)\\[5pt]

D01:
\((3,4)\)
&
\(\displaystyle -4(c_0+2c_2)q_0q_3\)\\[5pt]

D01:
\((4,9)\)
&
\(\displaystyle A q_0^2-Bq_3^2\)\\[5pt]

D01:
\((0,4)\)
&
\(\displaystyle -Bq_0^2+Aq_3^2\)\\[5pt]

D01:
\((5,5)\) \emph{(new in v6)}
&
\(\displaystyle -8c_1P^2\)\\

\end{longtable}

The entries related by \(1\leftrightarrow2\) are identified in the
executed record:
\(B_{44}=B_{77}\), \(B_{04}=B_{07}\),
\(B_{34}=B_{37}\), and \(B_{49}=B_{79}\).
This does not make \(B_{44}\) equal to \(B_{47}\) in D01 or D03.

\subsection{Tadpoles D02 and D12 and the \texorpdfstring{\((12)\)}{(12)} entries (new in v6)}

\textsf{[exec]} v4 did not print D02 and D12 as components.
The executed record gives them in full.
Put
\[
W=\tfrac12(c_0+2c_1+c_2),\qquad
W^\CR=\frac{(2\alpha-1)(4\alpha-1)}{6\alpha},\qquad
W^\Fr=\frac{2(\alpha-1)(2\alpha-1)}{3\alpha}.
\]
The only nonzero D02 entries, together with their transposes, are
\begin{align*}
\widehat B_{\mathrm{D02}}\big/\bigl(\tfrac{\ka^2}{\pi}W\bigr):\quad
&(01)(01),\ (02)(02):\ q_3^2;\qquad
(13)(13),\ (23)(23):\ q_0^2;\\
&(12)(12):\ P^2;\qquad (11)(22):\ -\tfrac12P^2;\\
&(01)(13),\ (02)(23):\ -q_0q_3;\qquad
(03)(11),\ (03)(22):\ q_0q_3;\\
&(00)(11),\ (00)(22):\ -\tfrac12q_3^2;\qquad
(11)(33),\ (22)(33):\ -\tfrac12q_0^2.
\end{align*}
This makes 23 nonzero entries.
They agree with the v3 tadpole comparison recorded in
\loc{CHECKPOINT_v4.md}.
Moreover
\[
\widehat B_{\mathrm{D12}}=0
\]
identically on both branches.

The table above omitted one key present in the executed record,
the \((12)(12)\) entry:
\[
\widehat B_{\mathrm{D01}}^{(12)(12)}
=\widehat B_{\mathrm{D03}}^{(12)(12)}
=-8c_1P^2\,\frac{\ka^2}{16\pi}
=-\frac{\ka^2(\alpha-1)(5\alpha-1)}{24\pi\alpha}P^2,
\]
on both branches.
It vanishes at \(\alpha=1\) and is continuous at \(\alpha=1/2\), so
Proposition~\ref{prop:polynomialchecks} extends to it.

\paragraph{Mixed helicity.}
\textsf{[exec]} The reduction attaches one coefficient \(c_S\) to a
monomial only when both lines sit in the same channel.
Mixed-channel monomials are logged and dropped.
They occur only on \emph{vector-type} keys, those with a single
transverse index on a leg:
\begin{itemize}
\item D01: \((01)(01)\), \((02)(02)\), \((13)(13)\), \((23)(23)\),
\((01)(13)\), \((02)(23)\), and their transposes;
\item D06: ghost \(11,22\);
\item D13: \((01)c^1\), \((02)c^2\), \((13)c^1\), \((23)c^2\).
\end{itemize}
No lattice key \(\{00,03,33,11,22,12\}\) has one.
The vector-type entries of D01, D06, and D13 are therefore
\emph{excluded} from the record and \emph{not} zero.
They need mixed-channel traces, which are \OPEN.

\subsection{Ghost and source-ghost sector}

Define
\begin{equation}
R=3c_0+4c_1,\qquad
T=4c_0+3c_1,\qquad
V=c_1-c_0.
\label{eq:RTV}
\end{equation}
Then the v4 ghost self-energy is
\begin{equation}
\widehat B_{\mathrm{D06}}
 =\frac{\ii\ka^2}{8\pi}R P^2
 \begin{pmatrix}1&0\\0&1\end{pmatrix},
\label{eq:D06compact}
\end{equation}
where row and column labels are \(0,3\).

The source-ghost entries are
\begin{equation}
\widehat B_{\mathrm{D13}}
 =\frac{\ii\ka^2}{16\pi}
 \begin{pmatrix}
 -2Tq_0 & 2Vq_3\\
 -Rq_3 & Rq_0\\
 -2Vq_0 & 2Tq_3
 \end{pmatrix},
\label{eq:D13compact}
\end{equation}
with tensor-source row labels \(0,3,9\) and ghost column labels
\(0,3\).

For explicit comparison,
\begin{align}
R&=\frac{(\alpha-1)(17\alpha-7)}{12\alpha},\nonumber\\
T&=\frac{(\alpha-1)(11\alpha-7)}{12\alpha},\qquad
V=\frac{\alpha-1}{2}.
\label{eq:RTVexpanded}
\end{align}
Equations~\eqref{eq:D06compact}--\eqref{eq:RTVexpanded} reproduce
all entries of the v4 ghost table without changing any factor
of \(\ii\), \(\pi\), or two.

The polynomials are branch independent because \(c_2\) does not
occur. This is an algebraic fact about the printed catalogue.
It does not prove that every domain manipulation of the corresponding
F and CR graphs is identical.

%======================================================================
\section{Exact checks available from the printed catalogue}
\label{sec:checks}
%======================================================================

\begin{proposition}[Algebraic preservation and endpoint checks]
\label{prop:polynomialchecks}
The coefficient tables above are equivalent to the displayed
rational expressions of the supplied v4 source.
Within their stated sector:
\begin{enumerate}[label=\textup{(\roman*)}]
\item every printed F entry vanishes at \(\alpha=1\);
\item every printed CR entry equals its F partner at
\(\alpha=1/2\);
\item all printed entries are continuous at
\(\alpha=1/2,2/3,4/5\);
\item D04+D05a has the sum \eqref{eq:phiExpanded};
\item the printed D06/D13 entries are independent of the
weight-two branch choice.
\end{enumerate}
These checks do not certify absent vertex arrays.
\end{proposition}

\begin{proof}
The endpoint claims follow because
\(c_0=c_1=c_2^\Fr=0\) at \(\alpha=1\), whereas
\(c_2^\Fr=c_2^\CR\) at \(\alpha=1/2\).
All formulas are rational with no denominator zero in the stated
sector.

For equivalence, substitute \eqref{eq:cvalues} into
Tables~\ref{tab:phiC} and \ref{tab:hhC}.
For instance,
\begin{align*}
c_0-4c_1+c_2^\CR
 &=\frac{2(\alpha^2+6\alpha-1)}{12\alpha},\\
c_0-4c_1+c_2^\Fr
 &=\frac{2(\alpha-1)(\alpha-5)}{12\alpha},\\
c_0+c_2^\CR
 &=\frac{2(11\alpha^2-6\alpha+1)}{12\alpha},\\
c_0+c_2^\Fr
 &=\frac{2(\alpha-1)(11\alpha-7)}{12\alpha},\\
A^\CR
 &=\frac{101\alpha^2-60\alpha+7}{12\alpha},\\
A^\Fr
 &=\frac{(\alpha-1)(101\alpha-55)}{12\alpha},\\
B&=\frac{11\alpha^2-12\alpha+1}{12\alpha}.
\end{align*}
The common factor \(\ka^2/(16\pi)\) gives the v4
D01 denominators and numerators.

Similarly,
\[
c_0+4c_1
 =\frac{(\alpha-1)(19\alpha-5)}{12\alpha},
\]
which checks the nontrivial D03 entry.
The scalar sum is immediate from Table~\ref{tab:phiC}.
Finally \eqref{eq:RTVexpanded} proves the ghost-table equivalence
and its branch independence.
\end{proof}

\subsection{A restricted branch-transversality check}

The following statement concerns the printed independent-coordinate
array, not an unproved identification with a geometric Hessian.

\begin{proposition}[Gauge annihilation of the printed D01 branch difference]
\label{prop:branchtransverse}
Order the coordinates as
\[
x=(h_{00},h_{03},h_{33},h_{11},h_{22}),
\qquad x_T=h_{11}+h_{22},
\]
and let
\[
u=(4q_3^2,-8q_0q_3,4q_0^2).
\]
With \(\nu=2-\alpha^{-1}\), the CR-minus-F difference in
Table~\ref{tab:hhC} is the bilinear form
\begin{equation}
\begin{split}
\Delta\widehat B_{\mathrm{D01}}(x,z)
 =\frac{\ka^2\nu}{16\pi}
 \bigl[
 &(u\cdot x_L)z_T+x_T(u\cdot z_L)\\
 &+P^2x_Tz_T
 \bigr],
\end{split}
\label{eq:branchbilinear}
\end{equation}
where \(x_L=(h_{00},h_{03},h_{33})\).
It annihilates the declared longitudinal gauge vector
\begin{equation}
k(c)=
(2q_0c_0,\ q_0c_3+q_3c_0,\ 2q_3c_3,\ 0,\ 0)
\label{eq:gaugevector}
\end{equation}
in either slot.
\end{proposition}

\begin{proof}
Only \(c_2\) changes, by \(\nu\).
Reading its coefficients from Table~\ref{tab:hhC} gives
\eqref{eq:branchbilinear}.
Since \(k_T=0\),
\[
u\cdot k_L
 =8q_0q_3^2c_0
 -8q_0q_3(q_0c_3+q_3c_0)
 +8q_0^2q_3c_3=0.
\]
Both slot contractions vanish.
\end{proof}

Independent-component weights matter here.
One must not apply an additional off-diagonal multiplicity to an
array already differentiated with respect to independent
\(h_{03}\) coordinates.
v6 adds the D02/D12 contribution in
Corollary~\ref{cor:hhcontent}.
In the complete \(hh\) sum the \(S_6,S_7\) parts of this difference
cancel against D02, which leaves the gauge-invariant term
\(-\nu\ka^2P^2f_0f_0'/(4\pi)\).

The mixed-helicity statement of v4 is now an executed check; its
precise form is given in Section~\ref{sec:catalogue}.

%======================================================================
\section{The ghost-equation row and its exact scope}
\label{sec:ghost}
%======================================================================

\subsection{An explicit longitudinal contraction}

For this algebraic check choose the covector
\[
p_\mu=(q_0,0,0,q_3),\qquad
p^\mu=(q_0,0,0,-q_3).
\]
This fixes the spatial momentum sign rather than leaving the
word ``upper'' ambiguous.

For the three displayed source rows of a symmetric tensor \(Y\),
define
\begin{align}
(\mathsf F_LY)_0
 &=\frac{q_0}{2}(Y_{00}+Y_{33})-q_3Y_{03},\nonumber\\
(\mathsf F_LY)_3
 &=q_0Y_{03}-\frac{q_3}{2}(Y_{00}+Y_{33}).
\label{eq:FL}
\end{align}
These are the de Donder contractions when the transverse trace
of \(Y\) vanishes.

For a general tensor, the full contraction
\[
(\mathsf FY)_\rho
 =p^\mu Y_{\mu\rho}
  -\frac12p_\rho\eta^{\mu\nu}Y_{\mu\nu}
\]
instead has the additional terms
\begin{equation}
(\mathsf FY)_0=(\mathsf F_LY)_0
 +\frac{q_0}{2}(Y_{11}+Y_{22}),\qquad
(\mathsf FY)_3=(\mathsf F_LY)_3
 +\frac{q_3}{2}(Y_{11}+Y_{22}).
\label{eq:fullF}
\end{equation}
The omitted transverse source trace is therefore a real
identification obligation, not a value that may be silently
set to zero.

\begin{theorem}[Exact projected ghost-equation identity]
\label{thm:ghost}
For the printed matrices \eqref{eq:D06compact} and
\eqref{eq:D13compact},
\begin{equation}
\widehat B_{\mathrm{D06}}
 =-2\eta_\parallel\,
 \mathsf F_L\widehat B_{\mathrm{D13}},
\qquad
\eta_\parallel=\diag(1,-1).
\label{eq:ghostrow}
\end{equation}
The equality is an exact polynomial identity for every
\(0<\alpha\le1\) for which \(c_0,c_1\) are given by
\eqref{eq:cvalues}. It is independent of \(c_2\).

To identify \eqref{eq:ghostrow} with the full graph
antighost equation additionally requires the complete source
normalization, the matched Fourier convention, and either a
zero transverse D13 trace or its explicitly included contribution
in \eqref{eq:fullF}.
\end{theorem}

\begin{proof}
The identity
\[
T+V=R
\]
follows from \eqref{eq:RTV}.
Applying \(\mathsf F_L\) to \eqref{eq:D13compact} gives
\[
\mathsf F_L\widehat B_{\mathrm{D13}}
 =\frac{\ii\ka^2}{16\pi}RP^2
 \begin{pmatrix}-1&0\\0&1\end{pmatrix}.
\]
Multiplication by \(-2\eta_\parallel\) yields
\eqref{eq:D06compact}. Formula~\eqref{eq:fullF} states
exactly the additional transverse term needed for the full
identification.
\end{proof}

This proves the two ratios reported in v4 on the printed
longitudinal projection:
\(-1/2\) in the \(00\) contraction and \(+1/2\) in the
\(33\) contraction, with zero off-diagonal entries.
It is not a proof that all source components or all 26
Ward-constrained combinations have been tested.

\begin{remark}[Transverse source trace (upgraded in v6)]
\label{rem:transtrace}
\textsf{[exec]} The executed D13 reduction covers all \(10\times4\)
source--ghost keys.
The rows \((11)\) and \((22)\) vanish for the ghost components
\(c^0,c^3\) on both branches.
On the executed matrix \(\mathsf F=\mathsf F_L\), so the additional
term in \eqref{eq:fullF} is computed to be zero, not assumed zero.
The ghost row therefore holds for the full de Donder contraction,
components \(\rho\in\{0,3\}\), on constant profiles.
The source normalization and the Fourier convention are those of
the archived code. Their identification with the XVI-P antighost
equation follows (C1)--(C3) of Section~\ref{sec:STdiag}.
\end{remark}

\subsection{The Maxwell--Friedrichs comparison}

The free photon inverse compatible with the scalar Friedrichs
ghost is the Maxwell--Friedrichs construction described in the
supplied I-B-G-S excerpt:
\[
C_1G^\pm_{(1)}=G^\pm_{(0)}C_1.
\]
With Friedrichs realizations in the vector and scalar ghost channels,
the supplied statement gives the Lorenz condition for conserved
compact sources and locally square-integrable packet field strength
\cite[{\loc{thm:spin1}}]{CheungIBGS}.

Residual gauge parameters there are scalar Friedrichs solutions,
in a stated solution or spectral-packet domain, with gradients
in the vector domain. They are not restricted to nonzero
spacetime-compact homogeneous solutions.

This is a free Maxwell comparison, not the interacting
diffeomorphism-ghost identity of D06/D13.
Its availability supplies neither the missing transverse source
rows nor a loop-level Slavnov--Taylor theorem.
\separation

%======================================================================
\section{The \texorpdfstring{\(hh\)}{hh} Slavnov--Taylor row: v4 diagnostic, diagnosis, and constant-profile closure}
\label{sec:ST}
%======================================================================

\subsection{The supplied attempted row (v4, historical)}

The v4 computation reports a test of
\begin{equation}
2B_{hh}(K_2c,h)
 +\lambda\,\Gamma_{Kc}(Eh,c)=0,
\label{eq:attemptedST}
\end{equation}
using
\[
B_{hh}=\mathrm{D01+D02+D03+D11+D12},
\]
a linearized Einstein operator \(E\), and a convention constant
\(\lambda\).

According to that record:
\begin{itemize}
\item no nonzero \(\alpha\)-independent convention choice gave
the desired CR/F closure;
\item the UNIT replacement \(c_S=1\) gave
\(B_{hh}(K_2c,h)=0\) but a nonzero source term;
\item 72 scanned conventions did not repair the test.
\end{itemize}
In v5 the code, arrays, and scan log were unavailable.
v6 archives them (\loc{code/st_test_const.py}, \loc{code/st_scan.py},
\loc{logs/v4/st_*.json}) and diagnoses the failure in
Section~\ref{sec:STdiag}.

\subsection{Why UNIT requires its own analytic control theorem}

The original inference that UNIT \emph{must} close merely because
it resembles the \(l_\perp=0\) flat integrand is too strong.
Restriction to a momentum slice need not preserve an identity
that holds only after integration in all momentum components.

\begin{example}[A full-space total derivative need not vanish on a slice]
\label{ex:slice}
Let \(F(l,k)\) be smooth and rapidly decreasing in all variables.
Then
\[
\int \dd l\,\dd^2k\,
 \partial_{k_1}\bigl(k_1F(l,k)\bigr)=0.
\]
But at \(k=0\),
\[
\left.
\partial_{k_1}\bigl(k_1F(l,k)\bigr)\right|_{k=0}
 =F(l,0),
\]
whose longitudinal integral need not vanish.
Thus full-momentum integration and transverse restriction do not
commute with this cancellation.
\end{example}

A valid UNIT control must therefore derive the intended row for
the actual reduced theory, source insertions, independent-component
normalization, regulator, and momentum-routing operations.
Longitudinal shift invariance alone does not establish that theorem.
Dimensional-reduction consistency is likewise not automatic
preservation of every proposed identity \cite{Stockinger}.

\begin{proposition}[No regularity verdict follows from the v4 test]
\label{prop:noSTverdict}
The supplied failure of \eqref{eq:attemptedST}, including its UNIT
diagnostic, proves neither failure nor validity of \(\Hreg\).
It does not test the complete set of 26 Ward-constrained combinations.
\end{proposition}

\begin{proof}
The complete row, its source and auxiliary-field terms, its
normalization, and a validated reduced control are not established
in the supplied test.
Several distinct missing choices can change its residual:
invariant versus gauge-fixed \(E\), auxiliary and antighost
contributions, source momentum routing, transverse source traces,
and independent-component multiplicities.

An unexplained residual of that unvalidated map cannot be identified
with the residual of the required full identity.
Conversely, failure of one candidate test does not establish
the required identity either.
\end{proof}

The ending of v5 for the \emph{v4 row} remains correct: that row
was \UNVALIDATED.
It was not ``cone anomaly found,'' and it was not ``BRST restored.''
v6 replaces it with a derived row, analysed below.

\subsection{Diagnosis of the v4 failure (new in v6)}
\label{sec:STdiag}

The frozen XVI-P antifield vertex has the closed form
\(K'_{(mn)}\,a\,(\partial_mc_n+\partial_nc_m)\), summed over the
independent coordinates \((mn)\), \(m\le n\) \cite{CheungP}.
Writing the BRST source term in the same independent coordinates
fixes three conventions that v4 left free:
\begin{enumerate}[label=\textup{(C\arabic*)}]
\item the ghost enters the vertex with an upper index, and
\(c_n=\eta_{nn}c^n\) appears in
\(\delta_ch_{(mn)}=p_mc_n+p_nc_m\);
\item the equation-of-motion factor is the derivative with respect
to the independent coordinate,
\[
\frac{\partial S}{\partial h_{(mn)}}=(2-\delta_{mn})E^{mn}(h),
\]
with \(E\) the linearized Einstein operator and indices raised by
\(\eta=\diag(1,-1,-1,-1)\);
\item consequently the \(K'_{(03)}\) coordinate carries weight
\(2E^{03}=-2E_{03}\).
In lower-index components this is weight \(w=-2\).
\end{enumerate}
\textsf{[exec]} The v4 convention scan (\loc{code/st_scan.py},
\loc{logs/v4/st_scan.json}) used only
positive weights \(w\in\{\tfrac12,1,2\}\).
It never contained \(w=-2\) on \((03)\) together with (C1).
This is why none of its 72 conventions closed.
The UNIT ``nonzero source term'' in v4 is produced by the same
weight error.

\subsection{An \emph{a priori} normalization}

v5 rightly warned that a fitted \(\lambda\) proves nothing.
In v6, \(\lambda\) is fixed before the loop data are used.
The XVI-P tree data give the antighost--ghost normalization
\(A_{\rm GH}=\tfrac12\).
They also give the quadratic-action normalization \(N=1\) of
\(E^{\rm up}\) on de Donder configurations.
\textsf{[exec]} \loc{code/st_norm.py} solves the de Donder
constraint for \((h_{00},h_{03})\), with \(h_{11}=h_{22}\) tied, and
returns
\[
\lambda_{\rm pred}=\frac{N}{\ii A_{\rm GH}}=-2\ii
\]
(\loc{logs/v6/st_norm.json}).

\subsection{Constant-profile closure on CR and F}

Let the axisymmetric lattice coordinates be
\(x=(h_{00},h_{03},h_{33},f_0)\), with
\(h_{11}=h_{22}=f_0\), and the ghost \(c=(c^0,c^3)\).
Put \(p_\mu=(q_0,0,0,q_3)\).
Define
\begin{equation}
\mathcal S_\lambda(c,h)=
\sum_{a,b}\widehat B^{hh}_{ab}\,(\delta_ch)_a\,h_b
+\lambda\sum_{(mn),\rho}
\widehat\Gamma^{\mathrm{D13}}_{(mn),\rho}\,c^\rho\,
(2-\delta_{mn})E^{mn}(h),
\label{eq:STv6}
\end{equation}
where
\(\widehat B^{hh}=\mathrm{D01+D02+D03+D11+D12}\)
denotes the executed constant-profile poles restricted to lattice
keys.

\begin{theorem}[Constant-profile \(hh/Kc\) Slavnov--Taylor row]
\label{thm:STclosure}
\textsf{[exec]}
For both CR and F and every \(1/2<\alpha<1\), the polynomial
\(\mathcal S_\lambda\) in
\((h,f_0,c,q_0,q_3)\) vanishes identically if and only if
\(\lambda=-2\ii=\lambda_{\rm pred}\).
More precisely:
\begin{enumerate}[label=\textup{(\roman*)}]
\item Sixteen monomials receive a nonzero graph contribution.
On twelve of them, the D01 and D03 contributions cancel among the
\(hh\) graphs.
D02, D11, and D12 contribute nothing to the row.
\item The remaining four monomials are
\(f_0c^0q_0^3,\ f_0c^0q_0q_3^2,\ f_0c^3q_0^2q_3,\ f_0c^3q_3^3\).
On them, the \(hh\) gauge variation has coefficients
\(\mp\ka^2(\alpha-1)/(4\pi)=\mp\ka^2V/(2\pi)\), and the antifield
term has coefficients
\(\pm\ii\ka^2(\alpha-1)/(8\pi)\).
Hence a single \(\alpha\)-independent \(\lambda=-2\ii\) cancels
all four.
\item The matter sub-row \(\mathrm{D11}+\mathrm{D12}\) is
transverse by itself.
\item The UNIT replacement \(c_S\equiv1\) makes both terms vanish
identically, a trivial closure.
So does the flat four-dimensional benchmark on the lattice
(\loc{code/flat_pole.py}).
Neither of these can determine \(\lambda\), which is why
\(\lambda\) is fixed \emph{a priori}.
\end{enumerate}
\end{theorem}

\begin{proof}[Evidence]
\loc{code/st_v6.py} forms \eqref{eq:STv6} symbolically from the
logged poles and extracts every monomial coefficient as an exact
rational function of \(\alpha\).
It solves for \(\lambda\) on the first equation and checks
\(\partial_\alpha\lambda=0\) and every residual \(=0\)
(\loc{logs/v6/st_v6_\{CR,F,UNIT,FLAT\}.json}).
The per-graph table is in \loc{logs/v6/st_v6\_stats.json}.

Section~\ref{sec:structures} gives the structural reason.
On the lattice, the gauge variation of the \(hh\) sum is
\(2(b_6+b_7)P^2f_0\,(p\cdot c)\), so
\(\mathcal S_\lambda\propto V\) on both branches.
The branch-dependent structure \(b_8\) is gauge invariant.

The validity of the flat tensor-averaging routine is checked
separately. The flat D11 covariant ratios \([6,1,-6,-1,8]\)
reproduce the frozen \texttt{P18\_D11} record
(\loc{logs/v6/flat\_vs\_p18.json}).
\end{proof}

\begin{remark}[Scope of the closure]
\label{rem:STscope}
Theorem~\ref{thm:STclosure} answers v5's slice objection
(Example~\ref{ex:slice}) for this row by direct computation, not
by appeal to shift invariance.
It is \emph{one} row, under these restrictions:
\begin{itemize}
\item the constant profile;
\item the axisymmetric lattice \(h_{11}=h_{22}\), without the
\((12)\) entry;
\item excluding the vector-type keys on which mixed-helicity
monomials were dropped (Section~\ref{sec:catalogue});
\item without the auxiliary-field terms that vanish on this
restriction;
\item and with the spectral reduction itself conditional on
Assumption~\ref{ass:reduction}.
\end{itemize}
It is not the 26 Ward-constrained combinations, and it is not a
full Slavnov--Taylor theorem on the cone.
Within that scope the honest status changes from \UNVALIDATED\
to \emph{closed}.
No anomaly is found, and no BRST restoration is claimed.
\end{remark}

%======================================================================
\section{Invariant decomposition of the constant-profile \texorpdfstring{\(hh\)}{hh} poles (new in v6)}
\label{sec:structures}
%======================================================================

On the axisymmetric lattice, write \(h\) for the longitudinal
\(2\times2\) block, \(\operatorname{tr}h=\eta^{ab}h_{ab}\),
\(h\cdot h'=h^{ab}h'_{ab}\), and
\(php=p^ap^bh_{ab}\) for \(a,b\in\{0,3\}\).
For two-dimensional symmetric matrices the Cayley--Hamilton
identity removes \(p\,hh'\,p\).
A symmetric \(SO(1,1)\)-invariant bilinear that is quadratic in
\(p\) is therefore a combination of six structures:
\begin{align*}
S_1&=P^2\operatorname{tr}h\operatorname{tr}h', &
S_2&=P^2\,h\cdot h', &
S_3&=(php)\operatorname{tr}h'+\operatorname{tr}h\,(ph'p),\\
S_6&=P^2(f_0\operatorname{tr}h'+\operatorname{tr}h\,f_0'), &
S_7&=(php)f_0'+f_0(ph'p), &
S_8&=P^2f_0f_0'.
\end{align*}

\begin{proposition}[Structure coefficients]
\label{prop:structures}
\textsf{[exec]}
Every executed constant-profile \(hh\) pole decomposes uniquely as
\(\sum_ib_iS_i\).
In units of \(\ka^2/\pi\), the coefficients are the following
branch-uniform linear forms in \((c_0,c_1,c_2)\):
\begin{center}
\small
\begin{tabular}{@{}lcccccc@{}}
\toprule
Graph & \(b_1\) & \(b_2\) & \(b_3\) & \(b_6\) & \(b_7\) & \(b_8\)\\
\midrule
D01 & \(-\tfrac1{16}c_0\) & \(\tfrac18c_0\) & \(0\) &
\(-\tfrac18c_0-\tfrac14c_1-\tfrac12c_2\) &
\(\tfrac14c_0+\tfrac12c_2\) &
\(\tfrac14c_0-\tfrac12c_1+\tfrac14c_2\)\\
D02 & \(0\) & \(0\) & \(0\) & \(W\) & \(-W\) & \(-W\)\\
D03 & \(\tfrac1{16}c_0\) & \(-\tfrac18c_0\) & \(0\) &
\(\tfrac38c_0\) & \(-\tfrac14c_0\) &
\(-\tfrac14c_0-\tfrac12c_1\)\\
D11 & \(0\) & \(0\) & \(0\) & \(0\) & \(0\) & \(\tfrac18c_0\)\\
D12 & \(0\) & \(0\) & \(0\) & \(0\) & \(0\) & \(0\)\\
\midrule
Sum & \(0\) & \(0\) & \(0\) &
\(\tfrac34(c_0+c_1)\) & \(-\tfrac12c_0-c_1\) &
\(-\tfrac38c_0-2c_1-\tfrac14c_2\)\\
\bottomrule
\end{tabular}
\end{center}
Here \(W=\tfrac12(c_0+2c_1+c_2)\).
The CR fit determines the forms uniquely.
On F, where \(c_2^\Fr=7c_0+6c_1\) makes the three coefficients
linearly dependent, the same forms reproduce every entry
(\loc{logs/v6/hh\_structures.json}).
UNIT (\(c\equiv1\)) gives the sum
\((0,0,0,\tfrac32,-\tfrac32,-\tfrac{21}8)\).
\end{proposition}

\begin{corollary}[Gauge and branch content of the complete \(hh\) sum]
\label{cor:hhcontent}
\leavevmode
\begin{enumerate}[label=\textup{(\roman*)}]
\item The purely longitudinal structures cancel:
\(b_1=b_2=b_3=0\).
This cancellation is between D01 and D03.
\item The only gauge-variant combination on the lattice is
\[
b_6+b_7=\tfrac14(c_0-c_1)\frac{\ka^2}{\pi}
=-\frac{\ka^2V}{4\pi}
=-\frac{\ka^2(\alpha-1)}{8\pi},
\]
which is branch independent.
This is exactly the term balanced by the antifield row in
Theorem~\ref{thm:STclosure}.
\item The CR--F difference of the complete \(hh\) sum is
\[
\Delta\widehat B^{hh}=-\frac{\nu\ka^2}{4\pi}P^2f_0f_0'.
\]
It is annihilated by the longitudinal gauge vector
\eqref{eq:gaugevector}, since \(k_T=0\).
The D01 difference of Proposition~\ref{prop:branchtransverse} has
\(S_6,S_7\) parts. They cancel against the D02 difference
\(\tfrac{\nu}{2}(S_6-S_7-S_8)\ka^2/\pi\).
This settles the v5 obligation that D02/D12 must be included in
the branch-transversality check.
\end{enumerate}
\end{corollary}

\begin{remark}[Relation to the 18-dimensional engine map]
Axisymmetric constant-profile data land in the six-dimensional
space above.
The programme's complete \(hh\) map is 18-dimensional
\cite{CheungConeV4,CheungConeV5}; it is computed by the
XVI-P engine in its own basis.
That basis also includes transverse-derivative and non-axisymmetric
structures, which constant profiles cannot see
(Proposition~\ref{prop:noninjective}).
The convention map that sends the six coefficients above to their
images in the engine's 18 coordinates is \(\OPEN\).
So, therefore, are the components of the 18-dimensional vector
outside the image.
They are not set to zero.
\end{remark}

%======================================================================
\section{What constant-profile data cannot determine}
\label{sec:missing}
%======================================================================

\subsection{The full basis projection}

The supplied programme counts a 34-dimensional apex basis, including
26 Ward-constrained combinations and six exact plus two
\(\phi\phi\) directions. Those counts are retained as supplied
basis metadata; the basis itself is not constructed from the
printed component table.

\begin{proposition}[Constant-profile restriction is not injective]
\label{prop:noninjective}
On any smooth external-jet class containing constant transverse
jets and quadratic transverse jets, restriction of apex-local
bilinear functionals to constant profiles need not be injective.
Consequently, constant-profile values alone do not determine
coefficients multiplying all transverse-derivative operators.
\end{proposition}

\begin{proof}
Consider the symmetric bilinear functional
\[
L(F,G)=
\int_{\WW}
\left(F\Delta_\perp G+G\Delta_\perp F\right)\,\dd^2y.
\]
It vanishes when both fields are transversely constant near the apex.
Choose \(F\) with constant transverse profile near zero and
\(G=r^2g(y)\) there, with a smooth outer cutoff.
Since \(\Delta_\perp r^2=4\) on the two-dimensional cone away
from the apex and the radial jet has that limit,
\(L(F,G)\) is nonzero for suitable longitudinal tests.
Thus the restriction has a nontrivial kernel.
\end{proof}

This elementary example does not declare the allowed physical
counterterm basis. It proves why the missing basis cannot be
recovered solely by evaluating the displayed restriction.

The single computed scalar combination is the longitudinal
\(\delta_{\WW}\phi\,\partial_\parallel^2\phi\)-type value in
\eqref{eq:phiExpanded}, in the supplied convention.
In v6 the second scalar combination, which belongs to a quadratic
transverse profile, is \emph{partial}.
Its D05a part is computed to be \(0\), and its D04 part is open
(Section~\ref{sec:weighted}).

The higher coefficient \(\mathcal D_b\) that v4 quoted without
derivation is now derived in Theorem~\ref{thm:weighted}.
The trace-cyclicity argument of v4 is not used: \(\mathcal D_b\)
comes from its own normalized weighted heat trace, and its domain
audit is the flip estimate of Proposition~\ref{prop:weightedflip}.

\subsection{Joint-apex condition \(\Jcond\)}

If Assumption~\ref{ass:reduction} holds, the constant-profile apex UV
pole is the longitudinal logarithm in
Theorem~\ref{thm:reduction}. This is the precise conditional
version of the spectral joint-apex assertion.

It is not an unconditional proof of \(\Jcond\) for all graphs
or all external jets.
General profiles require the corresponding weighted spectral
coefficients, derivative-domain control, and graph contractions.
For D01 the additional requirement
\((\mathrm H_{I,\mathrm{D01}})\) is now resolved in the Mellin
sense, and its prescription status is recorded in
Corollary~\ref{cor:HI}.

\subsection{Regularity at special cone openings}

The printed rational functions are continuous at
\[
\alpha\in\left\{\frac12,\frac23,\frac45\right\}.
\]
This is proved in Proposition~\ref{prop:polynomialchecks}.
Continuity of finitely many coefficients is not a uniform estimate
for the operator domains or nonconstant jets required by \(\Hreg\).

The v3 claim about
\(H^1_{(3)}=0\) modulo \(\Span\) is retained only as a
version-qualified supplied assertion.
The exact complex, span matrix, and proof are not present in v4,
so it is not used to infer the absence of a local obstruction.

\subsection{Evanescent data and finite linear certificates}

\paragraph{Evanescent coefficients (item (d)): partial.}
Lemma~\ref{lem:UV} shows that an \(O(\eps)\) dependence of the
angular continuation, or of the quasi-four-dimensional index
algebra, leaves the simple-pole coefficient unchanged.
So the four-dimensionally projected poles of
Section~\ref{sec:catalogue} are stable at pole level.
The spectral code follows the XVI-P DRED prescription, with
quasi-\(d\) loop momenta and quasi-four-dimensional vertex indices.

The evanescent \emph{coefficients} needed for finite parts and
later identities require the \(\tilde g\)-leg (evanescent-index)
vertex arrays.
These are not part of the frozen XVI-P exports used here.
They are \(\OPEN\).
No evanescent coefficient is replaced by its four-dimensional value.

\paragraph{Exact \(\Span\) certificate (item (f)): open.}
The XVI-P engine computes the counterterm span natively modulo a
prime.
An exact certificate requires either a port of the rank computation
to \(\mathbb Q(\ii)\) with fraction-free elimination, or a
specified nonzero minor evaluated exactly.
Neither an exact matrix export nor such a port was found in the
frozen trees as inspected.
Building one would mean modifying or extending the frozen engine,
which this revision does not do.
Memory is not the obstruction.
A modular rank, even at several primes, is evidence and not an
exact certificate, and it is not reported as one.

\paragraph{Remaining items.}
The following remain \(\OPEN\):
\begin{enumerate}[label=\textup{(\roman*)}]
\item the evanescent coefficients above;
\item the convention map of the six structure coefficients of
Section~\ref{sec:structures} into the 18-dimensional engine basis,
and the components outside its image;
\item the full 34-dimensional basis projection and the 26 Ward
combinations;
\item the exact \(\Span\) certificate;
\item the vector-type keys excluded in Section~\ref{sec:catalogue}.
\end{enumerate}
A numerical rank or one chosen minor does not certify a whole
unspecified contact family. The matrix, coefficient field,
normalizations, and completeness of its columns must first be
declared.

%======================================================================
\section{Companion comparisons and finite-report boundaries}
\label{sec:comparisons}
%======================================================================

\subsection{Smooth cores, one-point data, and infrared control}

The v4 intake reports two newly found registry identifiers:
\loc{I-B-VP-5-Ex1-B} and \loc{I-B-VP-5-Ex2-B}.
The historical count is relative to that reported intake snapshot,
not a claim about the present registry.

The supplied Ex1-B material concerns non-universality of complete
smooth-core finite parts and heat-normalization audits
\cite{CheungEx1B}. The supplied Ex2/Ex2-B material concerns
twisted scalar Friedrichs selection, a Dirac critical-window
selection, and operator-limit qualifications
\cite{CheungEx2,CheungEx2B}.
Neither supplies an interacting spin-2 CR graph calculation.

The supplied I-B-VP material separately reports compact-test
infrared control for the massless Friedrichs spectral covariance
and regular-host one-point formulas
\cite[{\loc{lem:infrared}, \loc{thm:friedrichs}}]{CheungVP}.
That infrared statement is not a proof that the present massless
loop integrals have no infrared poles.

The Ward-compatible extension freedom for one-point stress and
current distributions in I-B-VP-1 is another extension problem
\cite{CheungVP1}.
It does not replace the Maxwell--Friedrichs photon domain, choose
the spin-2 CR/F realization, or fix a two-point graph prescription.
Strong-resolvent convergence alone also does not imply convergence
of singular covariances or renormalized one-point observables.

\subsection{Continuum kernels and echo reports}

The free Green identities in I-B-G are tied to their spectral
test domain. They are not automatically an exact sequence on the
ordinary compact-spacetime test space
\cite[{\loc{thm:green}, \loc{rem:compact-exact}}]{CheungIBG}.

I-B-ECHO compares image-isolated windows, diffracted fronts,
short-window deficit sensitivity, and finite-precision neutral-core
reports \cite{CheungEcho}.
Those results do not identify the window functional with the apex
residue. Example~\ref{ex:image} already separates the heat constant
from a purely diffracted-front diagnostic.

If one later requests a finite report of a residue, three new
objects must be specified:
\[
\text{defined continuum coefficient}
\quad\longrightarrow\quad
\text{acquisition or supplied preparation}
\quad\longrightarrow\quad
\text{complete finite instrument}.
\]
Only the first is the target of the present loop calculation.
A finite code cannot create the missing coefficient or its Ward
identity.

\subsection{Finite Enter instruments and their ledgers}

For comparison, the finite destructive instrument has form
\[
\mathcal M(X)=
\sum_j\proj j_R\otimes D_jXD_j^\dagger,\qquad
D_j=\bra{\tau_j}_Q\otimes I_A,
\]
with
\[
\sum_jD_j^\dagger D_j=I,\qquad
\Tr_R\mathcal M(X)=\Tr_QX.
\]
These identities follow by matrix-unit expansion and remain true
with untouched finite references
\cite{CheungIB,CheungEn09}.
They consume a finite name and return a classical code with its
declared residual; they do not select the continuum domain.

For a supplied real signal \(s\), X-J-ST uses laws of the form
\[
P_s(F)=r,\qquad
P_s(\pm)=\frac{1-r}{2}
 \left(1\pm\tanh(\gamma s/\sigma)\right).
\]
For two supplied values,
\[
\TV(P_{s_1},P_{s_2})
 =\frac{1-r}{2}
 \abs{\tanh(\gamma s_1/\sigma)-\tanh(\gamma s_2/\sigma)}.
\]
Complete implementation allowances \(e_1,e_2\) reduce a guaranteed
ideal separation by at most \(e_1+e_2\), by the triangle inequality
\cite{CheungXJ,CheungEcho}.

Finite clock composition and resource fitting require the full
history, failures, controls, locations, and storage ledgers
\cite{CheungXK}. These facts price an already declared finite
interface. They neither acquire a continuum pole nor validate
a failed ST test.
Heat belongs to a separate physical implementation ledger
\cite{CheungIC,ReebWolf}.
\separation

%======================================================================
\section{Historical corrections and reproducibility obligations}
\label{sec:repro}
%======================================================================

\subsection{Corrections retained from v4}

The following v4 corrections to v3 remain:
\begin{enumerate}[label=\textup{(\roman*)}]
\item D06 is the \(\bar c c\) ghost self-energy, not
``the D01+D03 channel.''

\item The source reports confirmation of the D05a symmetry factor
\(1/2\) from Q02. Its displayed value is unchanged.

\item The closed constant-profile polynomials replace the earlier
unevaluated Mellin-residue status only in that restricted
reported computation. They do not supply nonconstant jets.

\item The registry log distinguishes the oct18c report of zero
new identifiers from the later oct18d report of two.
v6 rechecked it: 251 identifiers, unchanged.
\end{enumerate}

\subsection{Proof-level corrections made in v5 (retained)}

\begin{enumerate}[label=\textup{(\roman*)}]
\item The apex functional is explicitly bulk-subtracted.

\item Nonperiodic CR labels are justified through a singular
Bessel-channel replacement.

\item The flat CR remnant is identified as a one-sided limit.

\item The heat-trace formula and complete graph reduction have
separate hypotheses.

\item The projected ghost row has explicit momentum and
transverse-trace conventions.

\item The UNIT benchmark is not assumed valid merely from
longitudinal shift invariance.

\item The full \(hh\) tadpole arrays and higher-profile values are
not reconstructed from ``as v3'' (v6: now executed and printed).

\item Self-adjoint Schr\"odinger evolution, positive wave evolution,
and open-system access loss are distinguished in
Appendix~\ref{app:readout}.
\end{enumerate}

\subsection{Corrections in v6}

\begin{enumerate}[label=\textup{(\roman*)}]
\item The \(K'_{(03)}\) coordinate weight in the \(hh\) ST row is
\(-2\), not a positive weight (Section~\ref{sec:STdiag}).
\item The ST constant is \(\lambda=-2\ii\) in the normalization
\eqref{eq:STv6}.
The early-v6 checkpoint value \(-4\ii\) belongs to the v4
normalization, which has an extra factor 2 on the \(B\) term.
Both are recorded in \loc{CHECKPOINT_v6.md}.
\item The v4/v5 catalogue omitted the \((5,5)\) entries of D01 and
D03.
The executed record contains them, and they are added without
changing any printed entry.
\item ``Mixed-helicity cancellation'' is made precise: absent on
lattice keys, and excluded (not cancelled) on vector-type keys.
\item The early numerical \(\HI\) runs suffered catastrophic
cancellation in \(I_{-\nu}-I_\nu\).
They were replaced by the \(K_\nu\) form, analytic derivatives,
and an explicit small-\(r\) head.
Only the corrected run is reported.
\end{enumerate}

\subsection{Artifact inventory (v6)}

The executed record now consists of the following
(Appendix~\ref{app:artifacts}):
\begin{itemize}
\item the spectral reduction \loc{code/spectral_apex.py}, with the
graph rules copied read-only from XVI-P
\loc{amplitude_rules.py} and XVI-Q Q02;
\item per-graph logs in \loc{logs/v4/} for CR, F, and UNIT;
\item the v6 programs \loc{st_v6.py}, \loc{st_norm.py},
\loc{st_v6_stats.py}, \loc{flat_pole.py}, \loc{flat_vs_p18.py},
\loc{weighted_apex.py}, \loc{weighted_checks.py}, \loc{hi_d01.py},
\loc{hi_d01_exact.py}, \loc{hh_structures.py}, and
\loc{catalogue_check.py}, with their logs in \loc{logs/v6/};
\item full SHA-256 tree digests of XVI-P/Q/R, taken before and after
this revision.
\end{itemize}
The \loc{logs/v4/} spectral runs were executed on this machine
during the v4 session.
v6 re-reads and audits them, and reruns the derived checks.
v6 does \emph{not} rerun the spectral contractions themselves;
their per-graph timings are recorded in the JSON logs.
This replay-versus-rebuild distinction is recorded deliberately.

\subsection{Status of the computational-certificate obligations}

\begingroup\small
\begin{longtable}{@{}L{.44\linewidth}L{.50\linewidth}@{}}
\toprule
v5 obligation & v6 status\\
\midrule
\endfirsthead
\toprule
v5 obligation & v6 status\\
\midrule
\endhead
\bottomrule
\endfoot
1. Frozen source versions and complete hash manifests &
Done: full tree digests and full input-file digests.\\
2. Graph inventory, vertices, routing, phases, statistics, symmetry
factors, normalizations &
Done for D01--D06, D11--D13: archived in code, with values read
from frozen sources.\\
3. Helicity projectors and domain-changing operations &
Projectors archived. Domain: D01 composite exact
(Theorem~\ref{thm:HI}); other graphs per the v4 assessment.\\
4. Exact contracted spectral functions before the spectral
parameter is set to zero &
\OPEN: the code reduces at zero transverse momentum.\\
5. Regularized trace reductions and boundary defects &
Partial: Theorems~\ref{thm:apex}, \ref{thm:flip},
\ref{thm:weighted}, and \ref{thm:HI}.\\
6. UV/IR separation and regulator dependence &
Partial: Lemma~\ref{lem:UV}; (R4) open.\\
7. Full D02/D12 arrays and source transverse-trace rows &
Done (Section~\ref{sec:catalogue}).\\
8. Complete ST functional, derived control, raw residuals &
Done for one row within scope (Theorem~\ref{thm:STclosure}).
Full functional open.\\
9. Exact basis-projection and rank certificates & \OPEN.\\
10. Commands, versions, logs, failures, replay vs.\ rebuild &
Done (Appendix~\ref{app:artifacts}; corrections above).\\
\end{longtable}
\endgroup
Polynomial agreement and hash agreement alone still do not supply
items 4, 6, and 9.

%======================================================================
\section{Point-by-point response to the v5 audit (new in v6)}
\label{sec:response}
%======================================================================

The table answers every issue v5 raised against v4, in v5's order.
``Retained'' means v5's correction stands unchanged.
\textsf{[exec]} marks an executed v6 answer
(Appendix~\ref{app:artifacts}).

\begingroup
\small
\begin{longtable}{@{}L{.035\linewidth}L{.36\linewidth}L{.54\linewidth}@{}}
\toprule
\# & v5 issue against v4 & v6 answer\\
\midrule
\endfirsthead
\toprule
\# & v5 issue against v4 & v6 answer\\
\midrule
\endhead
\bottomrule
\endfoot
1 & Apex functional must be bulk-subtracted. &
Retained. Extended to the \(r^2\)-weighted trace
(Definition~\ref{def:weighted}, Theorem~\ref{thm:weighted}).\\
2 & Nonperiodic CR labels need a singular-channel justification. &
Retained (Theorem~\ref{thm:flip}). Weighted analogue added
(Proposition~\ref{prop:weightedflip}).\\
3 & Flat CR remnant is only a one-sided limit. &
Retained. New: the remnant is absent from the \(r^2\) moment
(Corollary~\ref{cor:Dvalues}).\\
4 & Heat trace and graph reduction need separate hypotheses. &
Retained (Assumption~\ref{ass:reduction}).
(R1) is discharged by the archived arrays.
(R3) is discharged for the D01 composite in the Mellin sense
(Theorem~\ref{thm:HI}).
(R2), (R4), and (R5) remain hypotheses.\\
5 & Catalogue identification conditional; vertex arrays and
execution artifacts not supplied. &
\textsf{[exec]} Arrays, routing, and prefactors are archived
(\loc{code/spectral_apex.py}).
All v5 entries match the logs exactly (\loc{code/catalogue_check.py}).
New entries: D01/D03 \((5,5)\), the full D02 array, and D12 \(=0\)
(Section~\ref{sec:catalogue}).\\
6 & Mixed-helicity cancellation is only a supplied report. &
\textsf{[exec]} No mixed-helicity monomial occurs on any lattice
key.
Such monomials occur only on vector-type keys, those with one
transverse index on a leg, in D01, D06, and D13 (listed in
Section~\ref{sec:catalogue}).
Those keys are excluded, not set to zero.\\
7 & D05a symmetry factor rests on a supplied Q02 comparison. &
\textsf{[exec]} Read directly from the regenerated
\loc{xviq_regen/Q02_bulk_graphs.json}: D05a \(1/2\), D12 \(1/2\),
D04 \(1\).
The engine prefactor is \(-\ka^2/2\), and the D05a value is
unchanged.\\
8 & Ghost row needs explicit momentum conventions and the
transverse source trace. &
Conventions retained (Theorem~\ref{thm:ghost}).
\textsf{[exec]} The D13 rows \((11),(22)\) are computed to be zero
for both ghost components. So \(\mathsf F=\mathsf F_L\) on the
executed matrix, and the transverse-trace obligation is discharged
for constant profiles.\\
9 & UNIT is not valid merely by longitudinal shift invariance
(slice example). &
Answered by direct computation instead of by that argument.
UNIT closes trivially, and \(\lambda\) is fixed a priori
(Theorem~\ref{thm:STclosure}).\\
10 & \(hh\) ST test is \UNVALIDATED; no verdict. &
Diagnosed: the \(K'_{(03)}\) weight should be \(-2\)
(Section~\ref{sec:STdiag}).
The derived row closes on CR and F with
\(\lambda=-2\ii=\lambda_{\rm pred}\), within the scope of
Remark~\ref{rem:STscope}.\\
11 & Full \(hh\) tadpole arrays not reconstructed from
``as v3''. &
\textsf{[exec]} D02 is printed in full and D12 \(=0\).
Both agree with the 24-coefficient v3 tadpole comparison recorded in
\loc{CHECKPOINT_v4.md}.\\
12 & D01 branch-transversality check omits D02/D12. &
Done: the complete \(hh\) branch difference is
\(-\nu\ka^2P^2f_0f_0'/(4\pi)\), which is gauge invariant
(Corollary~\ref{cor:hhcontent}).\\
13 & Independent-component weights must not be double counted. &
Confirmed and used. The \((2-\delta_{mn})\) weight is
\emph{derived}.
The flat D11 comparison with \texttt{P18\_D11} fixes off-diagonal
weight 2 in the covariant basis.\\
14 & 26 Ward combinations; 34-dimensional projection. &
\OPEN. One row is tested and closes.
The basis itself is still not constructed.\\
15 & Six exact plus two scalar basis values. &
One scalar value is retained.
The second is partial: D05a part \(=0\), D04 part \OPEN.
The six exact values are \OPEN.\\
16 & \(\mathcal D_b\) was quoted without derivation or value. &
Derived and checked (Section~\ref{sec:weighted}).\\
17 & \(\Jcond\) only conditional; \((\mathrm H_{I,\mathrm{D01}})\)
open. &
\((\mathrm H_{I,\mathrm{D01}})\): exact \(L^2\) obstruction and an
exactly zero Mellin value. The prescription choice is
\CONDITIONAL\ (Corollary~\ref{cor:HI}).
\(\Jcond\) for general profiles remains open.\\
18 & \(\Hreg\) at \(\alpha=1/2,2/3,4/5\). &
Coefficients remain continuous, and so do \(\mathcal D_b\) and
\(M(z)\) (no resonance for \(0<\nu<1\)).
The general statement is still open.\\
19 & \(H^1_{(3)}=0\) mod \(\Span\) only version-qualified. &
Unchanged; not used.\\
20 & Evanescent coefficients. &
Pole-level statement only (Section~\ref{sec:missing}).
The coefficients are \OPEN.\\
21 & Complete \(hh\) 18-dimensional map. &
Partial: the six-structure decomposition is complete.
The engine map is \OPEN.\\
22 & Exact \(\Span\) certificate. &
\OPEN, with reason (Section~\ref{sec:missing}).\\
23 & Truncated hashes cannot authenticate trees. &
Full tree digests before and after, and full file digests,
are in Appendix~\ref{app:artifacts}.\\
24 & Registry intake only historical. &
Quick recheck: 251 identifiers, identical in set and content to
oct18c\(\cup\)oct18d.
No change, so no \loc{new_oct18e/} directory was created.\\
25 & Ten-item reproducibility obligation list. &
Status per item in Section~\ref{sec:repro}.\\
26 & Compilation not asserted. &
Two pdfLaTeX passes with no unresolved references
(Appendix~\ref{app:declarations}).\\
27 & Readout analogy supplies no premise. &
Retained unchanged (Appendix~\ref{app:readout}).\\
\end{longtable}
\endgroup

%======================================================================
\section{Ledger and release}
\label{sec:ledger}
%======================================================================

\begingroup
\small
\begin{longtable}{@{}L{.50\linewidth}L{.44\linewidth}@{}}
\toprule
Item & Status in v6 (change from v5)\\
\midrule
\endfirsthead
\toprule
Item & Status in v6 (change from v5)\\
\midrule
\endhead
\bottomrule
\endfoot

Normalized twisted Friedrichs apex constant &
Proved: Theorem~\ref{thm:apex}.\\[3pt]

One exceptional regular-to-singular Bessel heat shift &
Proved: Theorem~\ref{thm:flip}.\\[3pt]

\(r^2\)-weighted apex coefficient \(\mathcal D_b\) and its shift &
\textbf{New}, proved: Theorem~\ref{thm:weighted} and
Proposition~\ref{prop:weightedflip}, with orbifold checks.\\[3pt]

Spectral reduction of a specified complete graph &
Conditional on Assumption~\ref{ass:reduction}.
(R1) discharged by archived arrays.\\[3pt]

D01--D06, D11--D13 constant-profile expressions &
\textbf{Upgraded}: executed and matching all v5 entries.
\((5,5)\), D02, and D12 added.
Graph identification still conditional on (R2)--(R5).\\[3pt]

D01 exceptional derivative tower \((\mathrm H_{I,\mathrm{D01}})\) &
\textbf{Upgraded}: exact \(L^2\) obstruction; Mellin value
\(=0\) exactly.
Prescription \CONDITIONAL.\\[3pt]

F flat limit and CR=F at \(\alpha=1/2\) &
Proved for the catalogue, including the new entries.\\[3pt]

Ghost-equation row &
Exact identity proved.
\textbf{Upgraded}: transverse trace computed to be zero.\\[3pt]

D02/D05a/D12 &
\textbf{Upgraded}: D02 printed, D12 \(=0\), D05a factor read from
Q02.\\[3pt]

Mixed-helicity terms &
\textbf{Upgraded}: absent on lattice keys; vector-type keys
excluded.\\[3pt]

\(hh\) Slavnov--Taylor row and UNIT &
\textbf{Upgraded} from \UNVALIDATED\ to \emph{closed in scope}:
CR and F with \(\lambda=-2\ii=\lambda_{\rm pred}\); UNIT trivially.\\[3pt]

Complete \(hh\) branch difference &
\textbf{New}: \(-\nu\ka^2P^2f_0f_0'/(4\pi)\), gauge invariant.\\[3pt]

26 Ward combinations; 34-dimensional projection & Not obtained.\\[3pt]

Six exact values & Not obtained.\\[3pt]

Second scalar value &
\textbf{Partial}: D05a part \(=0\); D04 part open.\\[3pt]

Joint-apex condition \(\Jcond\) &
Constant profile only, with the D01 composite now controlled.
General profiles open.\\[3pt]

\(\Hreg\) at \(1/2,2/3,4/5\) &
Coefficients continuous; general statement open.\\[3pt]

Evanescent coefficients &
\OPEN\ (pole-level statement only).\\[3pt]

Complete \(hh\) 18-dimensional map &
\textbf{Partial}: six-structure decomposition; engine map
open.\\[3pt]

Exact \(\Span\) certificate & \OPEN.\\[3pt]

Vector-type keys (one transverse index) &
\OPEN: they need mixed-channel traces.\\[3pt]

Registry intake & Rechecked; unchanged.\\[3pt]

Companion and catalogue connections &
Separations, not new premises or admission.\\[3pt]

\textbf{O.1 on the cone} & \textbf{\NOTCLOSED}.\\[3pt]

\loc{bulk_symmetry_restored} & \texttt{false}.\\
\end{longtable}
\endgroup

\begin{theorem}[Exact release of v6]
\label{thm:release}
Under the printed definitions and hypotheses:
\begin{enumerate}[label=\textup{(\roman*)}]
\item the normalized apex constants \(\Ccb_b\) and the weighted
coefficients \(\mathcal D_b\) are established, together with their
singular-channel shifts;
\item a complete graph satisfying the spectral reduction contract
has the longitudinal residue \eqref{eq:reduction};
\item the executed constant-profile poles equal the coefficient
catalogue of Section~\ref{sec:catalogue}, with the algebraic checks
proved above;
\item the projected ghost-equation row \eqref{eq:ghostrow} holds,
and the transverse source trace vanishes on the executed matrix;
\item the constant-profile \(hh/Kc\) Slavnov--Taylor row closes on
CR and F with the a-priori constant \(\lambda=-2\ii\);
\item the D01 composite has an exact \(L^2\) obstruction and zero
Mellin value;
\item the available data do not supply the full basis map, the six
exact values, the D04 quadratic-profile part, the evanescent
coefficients, the engine map, or an exact \(\Span\) certificate;
\item neither the companion separations nor the finite-readout
comparison changes the cone O.1 ending.
\end{enumerate}
No complete cone O.1 closure or bulk symmetry restoration is
released.
\end{theorem}

\begin{proof}
Theorems~\ref{thm:apex}, \ref{thm:flip}, and \ref{thm:weighted},
with Proposition~\ref{prop:weightedflip}, give (i).
Theorem~\ref{thm:reduction} gives (ii).
Proposition~\ref{prop:polynomialchecks} and the executed catalogue
check give (iii).
Theorem~\ref{thm:ghost} and the executed D13 rows give (iv).
Theorem~\ref{thm:STclosure} gives (v), and Theorem~\ref{thm:HI}
gives (vi).
Propositions~\ref{prop:noSTverdict} and \ref{prop:noninjective},
Remark~\ref{rem:STscope}, and Section~\ref{sec:missing} give (vii).

O.1 on the cone requires the full Ward/basis projection on the
prescribed cone, the evanescent and \(\Span\) data, and (R2)--(R5)
for every graph.
These are listed as open, and none is supplied by (i)--(vi).
The cone and CR branch remain inputs.
This gives (viii) and the release boundary.
\end{proof}

\scopebox{Conclusion}{%
v6 closes what the constant-profile record can close:
\begin{itemize}
\item the catalogue matches its executed source;
\item the v4 Slavnov--Taylor failure was a convention error, and
the derived row closes with an a-priori constant;
\item the D01 obstruction is exact, and it is Mellin-null;
\item a second apex moment is proved.
\end{itemize}
A closed row is not the complete Slavnov--Taylor identity.
A constant profile is not the full counterterm basis.
A Mellin value is not a prescription theorem.
The cone and CR domain remain inputs.
\[
\boxed{
\text{O.1 on the cone}=\NOTCLOSED,\qquad
\texttt{bulk\_symmetry\_restored}=\texttt{false}.
}
\]
}

%======================================================================
\appendix
\section{Self-adjoint apex dynamics and the readout analogy}
\label{app:readout}
%======================================================================

\subsection{Boundary form and flux}

For a real radial Bessel expression
\[
L_\nu=-\partial_r^2-r^{-1}\partial_r+\nu^2r^{-2}
\]
in \(L^2(r\,\dd r)\), integration by parts gives the endpoint
boundary pairing, up to the chosen outward orientation,
\begin{equation}
\mathfrak b(u,v)=
\lim_{r\downarrow0}
 r\left(\overline u\,v'
       -\overline{u'}\,v\right).
\label{eq:boundaryform}
\end{equation}
A self-adjoint domain makes this pairing vanish on every pair
of domain vectors, with the appropriate sum if channels are
coupled.

\begin{proposition}[No absorbing apex in the declared closed realization]
\label{prop:flux}
Let \(A\) be a self-adjoint realization with the declared closed
system boundary conditions.
Then:
\begin{enumerate}[label=\textup{(\roman*)}]
\item \(e^{-\ii tA}\) is unitary and preserves Hilbert norm;
\item the integrated endpoint probability flux of sufficiently
regular domain states vanishes;
\item a normalizable bound state is a state of that realization,
not an absorbing output channel.
\end{enumerate}
\end{proposition}

\begin{proof}
The spectral theorem gives a unitary group \(e^{-\ii tA}\).
Equivalently, for differentiable domain solutions,
\[
\frac{\dd}{\dd t}\norm{u(t)}^2
 =2\operatorname{Re}\pair{u(t)}{-\ii Au(t)}=0.
\]
The integrated radial current is proportional to
\(\operatorname{Im}(r\overline u\,u')\).
The vanishing of \(\mathfrak b(u,u)\) gives zero net boundary
flux. A bound eigenvector evolves by a phase
\(e^{-\ii t\lambda}\); its norm is not lost.
\end{proof}

This is standard self-adjoint-extension theory
\cite{KayStuder,DerezinskiRichard}.
For a channel-coupled boundary law, the proposition concerns the
total flux, not necessarily each component separately.

\begin{remark}[A negative bound state and wave stability]
Self-adjointness alone does not imply that
\(\partial_t^2u+Au=0\) is a stable positive-energy wave theory.
If \(Au=-\sigma^2u\), its wave component solves
\[
\ddot a-\sigma^2a=0.
\]
The sine quotient becomes \(\sinh(\sigma t)/\sigma\).
One must not call this positive-frequency evolution
\(e^{-\ii t\sqrt A}\) without a nonnegative operator.

The F and stated scale-invariant CR Bessel realizations above
are nonnegative in their declared range.
A different non-Friedrichs extension with a negative bound state
retains self-adjoint Schr\"odinger evolution but requires this
separate wave-stability qualification.
\end{remark}

Ghost and gauge-fixed tensor fields also do not automatically
carry a positive probability interpretation on physical
cohomology. No such identification is inferred from the radial
boundary calculation.

\subsection{What a loss channel would additionally require}

For the same complete single-particle realization, an absorbing
endpoint would require a different, dissipative boundary model
rather than the self-adjoint one just analyzed.
If the time evolution is written as \(e^{-tG}\), its contraction
generator is expressed by the appropriate maximal-accretive
condition on \(G\); the sign convention should not be transferred
uncritically to a Hamiltonian written in another convention.

There is another possibility: couple the cone to additional quantum
degrees of freedom and then restrict access to a subsystem.
The complete enlarged dynamics may remain self-adjoint and unitary,
while its accessible reduced channel is lossy.
Thus the unrestricted statement
``every form of loss requires a non-self-adjoint Hamiltonian''
would be false.

Neither an absorbing boundary condition nor an apex reservoir
coupling is constructed by this paper.
Both would change its declared model and require a new
gauge/BRST analysis.

\subsection{The precise I-C comparison}

I-C distinguishes destructive finite readout from coherent return
and from thermodynamic implementation \cite{CheungIC}.
Its six statements are relevant only as the following separations.

\begin{description}[style=nextline,leftmargin=0pt]
\item[A: recovery metric.]
If the complete accessible output of an unknown \(d\)-dimensional
code is classical, with no additional message-bearing quantum
resource available to the decoder, every deterministic decoder
\(\mathcal D\) obeys
\[
F_e(\mathcal D\mathcal C)\le\frac1d,\qquad
\frac12\norm{\mathcal D\mathcal C-\id}_\diamond
 \ge1-\frac1d.
\]
The distinct Haar-average pure-state infidelity is at least
\((d-1)/(d+1)\), using
\(F_{\rm av}=(dF_e+1)/(d+1)\)
\cite{HSR,NielsenAverage}.

\item[B: append-only entropy.]
For one unconditioned retained history,
\[
H(X_1,\ldots,X_{n+1})
 =H(X_1,\ldots,X_n)
  +H(X_{n+1}\mid X_1,\ldots,X_n).
\]
This monotonicity is not a Page curve for unmeasured quantum
radiation.

\item[C: conditional thermal accounting.]
For an initially independent Gibbs bath and a complete unitary
step,
\[
\beta_{\rm th}Q_B
 =S_{\rm n}(\rho_M)-S_{\rm n}(\omega_M)
  +I_{\rm n}(M:B)_\omega
  +D_{\rm n}(\omega_B\Vert\tau_B).
\]
A strictly positive complete non-bath entropy decrease implies
positive heat. The equality itself does not require that decrease,
and a readout channel does not fix the heat bill
\cite{ReebWolf}.

\item[D: complete classical exterior.]
Jointly classical heat labels and finite postprocessing do not
evade the classical-cut recovery obstruction. Quantum access to
the bath is a different interface.

\item[E: endpoint.]
A finite destructive prefix does not select physical evaporation,
a remnant, or a last-moment change to a quantum output.

\item[F: classical persistence.]
Persistent observables commute when the retained accessible
output is declared classical-only. Returned quantum residuals
and prior transfers cannot be omitted from the complete account.
\end{description}

For completeness, the recovery obstruction follows from the
measure-and-prepare form
\[
\mathcal D\mathcal C(\rho)
 =\sum_x\Tr(E_x\rho)\sigma_x.
\]
Its referenced output is separable.
For the maximally entangled state \(\Phi_d\), every product state
has overlap at most \(1/d\), since
\[
\Tr[\Phi_d(a\otimes b)]
 =\frac1d\Tr(a^{\mathsf T}b)\le\frac1d.
\]
Convexity and testing the Bell projector give the two displayed
bounds. This argument concerns a finite accessible channel,
not an apex extension.

\subsection{Why the analogy supplies no O.1 premise}

No entanglement-breaking map, classical-only radiation algebra,
bath-reset process, or finite acquisition instrument is part of
the cone residue calculation.
Conversely, the self-adjoint apex realizations do not choose such
an interface.

The correct conclusion is therefore:
\[
\boxed{
\begin{gathered}
\text{I-C supplies no usable analytical input for cone O.1;}\\
\text{self-adjoint endpoint data are not destructive readout;}\\
\text{no Ward identity or apex pole follows from the analogy.}
\end{gathered}}
\]
This is a separation, not a premise of any graph calculation
and not a claim that reduced access loss is impossible in a
different enlarged model.

%======================================================================
\section{Source and status declarations}
\label{app:declarations}
%======================================================================

\paragraph{Analytical inputs.}
The Bessel heat representation, standard Bessel transform identities,
self-adjoint functional calculus, and conventional dimensional
regularization background are identified analytical inputs.
The apex normalization, channel-shift integral, conditional reduction
argument, and polynomial identities are proved explicitly.

\paragraph{Companion source scope.}
The supplied v4 source, supplied I-C sources, version-qualified
excerpts, and programme descriptions are distinguished from
complete independently inspected manuscripts.
Companion manuscripts were not re-read for v6; their citations
carry the mark \(\dagger\).
No inaccessible source is claimed to have been read in full.

\paragraph{Primary literature.}
The bibliography identifies primary analytical and quantum-channel
sources for attribution and normalization.
The original shorthand citations ``FM97'' and ``S11'' are not used
as substitutes for precise hypotheses.
The normalized apex constant is proved here independently of those
abbreviations.

\paragraph{Execution records (v6).}
Statements tagged \textsf{[exec]} were executed on the shared
working machine and are logged.
Software: Python~3.13.5, SymPy~1.14.0, mpmath~1.3.0, and
pdfTeX~3.141592653-2.6-1.40.26 (TeX Live 2025).
The XVI-P/Q/R trees were hashed before and after this revision with
the command of Appendix~\ref{app:artifacts}; the digests are
identical.
No file in those trees was edited.
Nothing was pushed or written to any site.
The registry recheck read the public listing only.
The v4-era statements that v6 did not rerun stay attributed to
the supplied record.

\paragraph{Compilation and review.}
This source was compiled with two pdfLaTeX passes. It has an
internal bibliography, no external figure, data, or bibliography
file, and no unresolved references.
No proof-assistant check, independent referee review, physical
implementation, or journal acceptance is asserted.
In particular, \textsf{[exec]} is not independent review.

\paragraph{Submission obligation.}
A submission making the full loop-coefficient claim must attach or
archive the missing graph certificate and obtain the author's
approval of the mathematical scope and source versions.
The present conditional theorem must not be relabeled an
independently reproduced complete computation; \textsf{[exec]} means
reproducible on archived inputs, not independently reviewed.

\paragraph{Unchanged non-claims.}
The cone and CR extension are prescribed inputs.
Companion connections and OpenAI-attributed catalogue records
are documentary separations only.
No catalogue proof is adopted or independently verified.
O.1 on the cone stays \(\NOTCLOSED\);
\loc{bulk_symmetry_restored} stays \texttt{false}.
XVI-N is not discharged, the downstream XV-B admission exit is
not opened, and the H2/I3/J2 analytical chain is not completed.
No Einstein dynamics, physical tensor-core identification,
spacetime emergence, or inherited characteristic quantum
coefficient is constructed.

%======================================================================
\section{Artifact manifest and rerun commands (new in v6)}
\label{app:artifacts}
%======================================================================

Paths are relative to \loc{/workspace/XVI-O1-CONE/}.
\texttt{RO:} denotes the read-only copies under
\loc{/tmp/o1cone_ro/}; their digests equal those of the originals in
the frozen trees.
The frozen-tree digest is
\begin{center}\footnotesize
\verb!cd /workspace/XVI-$t && find . -type f -not -path './.pytest_cache/*'!\\
\verb!  -not -name '*.pyc' | sort | xargs sha256sum | sha256sum!
\end{center}

\begingroup\small
\begin{longtable}{@{}L{.10\linewidth}L{.52\linewidth}L{.25\linewidth}@{}}
\toprule
Tree & SHA-256 (before and after v6) & before vs.\ after\\
\midrule
\endfirsthead
\bottomrule
\endfoot
XVI-P & \texttt{\scriptsize ec965c15ae25cd3bb9cc2b3a134e966b}\newline\texttt{\scriptsize 19a1803c78d1fe5bf7554da2a8b3aaa6} & identical\\
XVI-Q & \texttt{\scriptsize 0df9a0d10409d31ac0c2d0a47197e3f6}\newline\texttt{\scriptsize e03c7ae258ec5fb775137f66cb0493c4} & identical\\
XVI-R & \texttt{\scriptsize 57187bc6825f5c6ae88633b7279c6c0b}\newline\texttt{\scriptsize e57f7c4918fdbc525da4aaf958331edb} & identical\\
\end{longtable}
\endgroup

\begingroup\small
\begin{longtable}{@{}L{.50\linewidth}L{.44\linewidth}@{}}
\toprule
File & SHA-256\\
\midrule
\endfirsthead
\toprule
File & SHA-256\\
\midrule
\endhead
\bottomrule
\endfoot
\multicolumn{2}{@{}l}{\textit{Base source (read-only copy)}}\\
{\footnotesize\loc{XVI-O1-CONE-v5.tex}} & \texttt{\scriptsize e00e49f548bdc3dd1da94411e7438f73}\newline\texttt{\scriptsize c873aa3e2b3d24bae3d74f2bd3a6985b}\\
\multicolumn{2}{@{}l}{\textit{Programs}}\\
{\footnotesize\loc{code/spectral_apex.py}} & \texttt{\scriptsize cb22811e721c932766a576cb15626015}\newline\texttt{\scriptsize 584f9cbc3cc3696f793e664c43debde2}\\
{\footnotesize\loc{code/st_test_const.py}} & \texttt{\scriptsize dd87dfd55dbf1a730ad31d728fced6d9}\newline\texttt{\scriptsize aa1f35c59afc0b5ab0f8072f9b81f0ee}\\
{\footnotesize\loc{code/st_scan.py}} & \texttt{\scriptsize 4a03c966a77b60a028e57163aef7f264}\newline\texttt{\scriptsize 5118c797ec3d55fc5c44ac8ee63cf759}\\
{\footnotesize\loc{code/catalogue_check.py}} & \texttt{\scriptsize 851f295b79cf2a483e6f9de4558d0f84}\newline\texttt{\scriptsize 2082bd7eff46efd1866d2ad93c3180df}\\
{\footnotesize\loc{code/st_v6.py}} & \texttt{\scriptsize 5dba49edc69ff70ca01cb0a428aef450}\newline\texttt{\scriptsize 262c117f19a6d86074acf4dc89cb6155}\\
{\footnotesize\loc{code/st_norm.py}} & \texttt{\scriptsize da1c19a9f75926f26490d44129fc21fa}\newline\texttt{\scriptsize 952b97384d4784a7c79219e5f1ff0f31}\\
{\footnotesize\loc{code/st_v6_stats.py}} & \texttt{\scriptsize 36ec924f70166b004720d547c1bf6cac}\newline\texttt{\scriptsize 4ae2283aec8321ea7db2234639241c31}\\
{\footnotesize\loc{code/flat_pole.py}} & \texttt{\scriptsize 7d77f198d4cd6874c6a88f8bff554835}\newline\texttt{\scriptsize 97b22ed74be4dc66bc7a2bbfd9423dc2}\\
{\footnotesize\loc{code/flat_vs_p18.py}} & \texttt{\scriptsize 4bee961081e64fa4ca7de67baa92bd51}\newline\texttt{\scriptsize 5bbd3307acaf317390c9fa5c68912561}\\
{\footnotesize\loc{code/weighted_apex.py}} & \texttt{\scriptsize a8a3ffd0742a5fdc5d176d8b259d4701}\newline\texttt{\scriptsize 1051513ff48191fcc22231e94cb2eada}\\
{\footnotesize\loc{code/weighted_checks.py}} & \texttt{\scriptsize a14d14e9112ab3603c4d77bbaed10cf1}\newline\texttt{\scriptsize dc70e67357ac4678d116de89eb317834}\\
{\footnotesize\loc{code/hi_d01.py}} & \texttt{\scriptsize f3e7ee7f82bde0c6aa85d4fc57e2807a}\newline\texttt{\scriptsize 47fe95a53b9af5348d31b206a0e524d5}\\
{\footnotesize\loc{code/hi_d01_exact.py}} & \texttt{\scriptsize c4e7d0271a0890d604e0ae1e8b300bb5}\newline\texttt{\scriptsize 3bdaa64402e1188faf72247133d0f118}\\
{\footnotesize\loc{code/hh_structures.py}} & \texttt{\scriptsize adc0bdbe0482e92d59bd26de05b90c41}\newline\texttt{\scriptsize e5d90e7dcb390c6ef6440f9d3c3f3f23}\\
{\footnotesize\loc{code/build_v6.py}} & \texttt{\scriptsize 40cd45f3b622260020cfb6f52a91904f}\newline\texttt{\scriptsize c2894a457be599fb12faec260b09663e}\\
{\footnotesize\loc{code/artifacts_v6.py}} & \texttt{\scriptsize d602a6ab9b12105f02baa274c2c212c3}\newline\texttt{\scriptsize 624c1c30d96c490cc8bb95e5e9020d00}\\
\multicolumn{2}{@{}l}{\textit{Spectral logs (v4 session, audited in v6)}}\\
{\footnotesize\loc{logs/v4/spec_D01_CR.json}} & \texttt{\scriptsize ee0c4766a8282d4969366eb44070617f}\newline\texttt{\scriptsize 715b535fc72c3760eb47e80130419a3e}\\
{\footnotesize\loc{logs/v4/spec_D01_F.json}} & \texttt{\scriptsize 7ea9fe6077c7b28d1bbad1f25fa7325a}\newline\texttt{\scriptsize a88e377227f87683d3a3ee570b9da819}\\
{\footnotesize\loc{logs/v4/spec_D01_UNIT.json}} & \texttt{\scriptsize 0923e4c51b68021f6fa77d734cdf25ec}\newline\texttt{\scriptsize fe758385a542b26207f1c97819b686a6}\\
{\footnotesize\loc{logs/v4/spec_D02_CR.json}} & \texttt{\scriptsize 679ac21d06a9b8de0657f15fbcb116cd}\newline\texttt{\scriptsize 6fbfcf5af7a5e2e56abd99cf2bde6398}\\
{\footnotesize\loc{logs/v4/spec_D02_F.json}} & \texttt{\scriptsize 827cc09f0945b69a9d0cba1bacf523df}\newline\texttt{\scriptsize 9e08424743220728034182a364de5d65}\\
{\footnotesize\loc{logs/v4/spec_D02_UNIT.json}} & \texttt{\scriptsize 2f12785331b61aa9534515191e916f46}\newline\texttt{\scriptsize cbba7d418c965b1e221140ab2bb36ee3}\\
{\footnotesize\loc{logs/v4/spec_D03_CR.json}} & \texttt{\scriptsize 6cae7e5ce69623b004ca177af9c4971b}\newline\texttt{\scriptsize 469976b704638810f978e68911a5e78f}\\
{\footnotesize\loc{logs/v4/spec_D03_F.json}} & \texttt{\scriptsize 7545563c038b73d232990a7b636055c1}\newline\texttt{\scriptsize 4687293c67e3cad76ba9fdd703261ef9}\\
{\footnotesize\loc{logs/v4/spec_D03_UNIT.json}} & \texttt{\scriptsize 8aee1bbdc1f7e6c31a378ae418054220}\newline\texttt{\scriptsize e605fd04314b9fc26e81b09ea34fa415}\\
{\footnotesize\loc{logs/v4/spec_D04_CR.json}} & \texttt{\scriptsize af01b84ec1891fff4c95d0ff43fa0854}\newline\texttt{\scriptsize 17ceaff7e876395026f3a506b84cbd72}\\
{\footnotesize\loc{logs/v4/spec_D04_F.json}} & \texttt{\scriptsize 7a209c8c669ca21a8119e168cfb40e62}\newline\texttt{\scriptsize e0b9642939dd9ef27a7307b07e6f4a2b}\\
{\footnotesize\loc{logs/v4/spec_D05a_CR.json}} & \texttt{\scriptsize c23d25b994b17fde9aa78a63bf520927}\newline\texttt{\scriptsize 33d452568a10ab7af77f828e4ddcbe03}\\
{\footnotesize\loc{logs/v4/spec_D05a_F.json}} & \texttt{\scriptsize c56f0489caae6064ba22ee9de188a66e}\newline\texttt{\scriptsize 4b173b11648e1a051c62fb387a5f3d9e}\\
{\footnotesize\loc{logs/v4/spec_D06_CR.json}} & \texttt{\scriptsize d759c4b910babacbc60b2aa2746da136}\newline\texttt{\scriptsize 669b2d40bbc3952e48c15bdd10f2dbcd}\\
{\footnotesize\loc{logs/v4/spec_D06_F.json}} & \texttt{\scriptsize 2fead6732f276e532fedc4d0aaf8bd27}\newline\texttt{\scriptsize 19b584ff7e5f23fa82701cef29c45652}\\
{\footnotesize\loc{logs/v4/spec_D06_UNIT.json}} & \texttt{\scriptsize a75b8ab0791896a4d64c06d14724c7f1}\newline\texttt{\scriptsize dccb5e0eed5358a6c40633b2190f360c}\\
{\footnotesize\loc{logs/v4/spec_D11_CR.json}} & \texttt{\scriptsize 6afa59f62437354f5dff58693c421e4b}\newline\texttt{\scriptsize 53867224bdd11b4b5ea47303f79be971}\\
{\footnotesize\loc{logs/v4/spec_D11_F.json}} & \texttt{\scriptsize b30291938bc5878556b69e6a2f8f4b88}\newline\texttt{\scriptsize 5ca803e5f3da1c0a60bb79727f1f55bb}\\
{\footnotesize\loc{logs/v4/spec_D11_UNIT.json}} & \texttt{\scriptsize f099bb41ff2ccbcfeb0bd568f3daf002}\newline\texttt{\scriptsize ec164a91d63ccf216b6aaf9347306508}\\
{\footnotesize\loc{logs/v4/spec_D12_CR.json}} & \texttt{\scriptsize d5f4ef51fd1adcf1e163422b76dc3421}\newline\texttt{\scriptsize 2a34a0c285cae006cfdda5fa1c2cc948}\\
{\footnotesize\loc{logs/v4/spec_D12_F.json}} & \texttt{\scriptsize 59937795d9674227954d36157ffb2d28}\newline\texttt{\scriptsize 08d38244e9346c55f416103b01c4abcc}\\
{\footnotesize\loc{logs/v4/spec_D12_UNIT.json}} & \texttt{\scriptsize b92e6a75b387400af13d4357cd5039dc}\newline\texttt{\scriptsize 5ce58b98f3a7e0c531048d78f39be508}\\
{\footnotesize\loc{logs/v4/spec_D13_CR.json}} & \texttt{\scriptsize 322b2653b807c47465cf9972da6acf5e}\newline\texttt{\scriptsize a0ef566b69374dfecb86ff402c497552}\\
{\footnotesize\loc{logs/v4/spec_D13_F.json}} & \texttt{\scriptsize ef4ef5e3cef8f2e9339a82e4f9f9d823}\newline\texttt{\scriptsize 844a13c2ebb10256bed893c6fb4b8f15}\\
{\footnotesize\loc{logs/v4/spec_D13_UNIT.json}} & \texttt{\scriptsize 9670e118cdd948d839d714242ba142af}\newline\texttt{\scriptsize e30cc1bf994c985d41ab2e8b583935a8}\\
{\footnotesize\loc{logs/v4/st_const_CR.json}} & \texttt{\scriptsize 23897227e71dea9022e56edf2441ac42}\newline\texttt{\scriptsize 56cbd97149036b3e46dccd3521ed64ff}\\
{\footnotesize\loc{logs/v4/st_const_F.json}} & \texttt{\scriptsize 56ca503cb7b1dd02db11bbc18658a373}\newline\texttt{\scriptsize dabeaaa579e7b39baf3b8ebb9e4e680d}\\
{\footnotesize\loc{logs/v4/st_scan.json}} & \texttt{\scriptsize e86875743f6e5c99852456e3d4b49293}\newline\texttt{\scriptsize 0727f4e0724a4f7dd1111e2ae41628db}\\
\multicolumn{2}{@{}l}{\textit{v6 logs}}\\
{\footnotesize\loc{logs/v6/catalogue_check.json}} & \texttt{\scriptsize 1a436314aaf0c5a7ca7f8cc7a7a09be0}\newline\texttt{\scriptsize a30394d80a1bbe2023fae6cca35f3e53}\\
{\footnotesize\loc{logs/v6/flat_D01.json}} & \texttt{\scriptsize 29293d0cf86d9a9091df2b6c5c613ea9}\newline\texttt{\scriptsize df2b501563164eca1ad4bd0db62dc86d}\\
{\footnotesize\loc{logs/v6/flat_D02.json}} & \texttt{\scriptsize 80804ff73c1b02d05f0b720df14b957d}\newline\texttt{\scriptsize fb73a38e0a9adeaa717e532bc1fcdf84}\\
{\footnotesize\loc{logs/v6/flat_D03.json}} & \texttt{\scriptsize ff09567fc9a68e73aa0b5c963359b111}\newline\texttt{\scriptsize 8445d807a9707e37c8b087b4dd0c31a9}\\
{\footnotesize\loc{logs/v6/flat_D06.json}} & \texttt{\scriptsize 346f4bf0cb4072a5f8cd38e2c725c069}\newline\texttt{\scriptsize 63bce2f8714ca50d7d92b7bae81d7332}\\
{\footnotesize\loc{logs/v6/flat_D11.json}} & \texttt{\scriptsize 9137e362feca3054545742d7cea4e439}\newline\texttt{\scriptsize 9fecddd117bdce0093bb79cb3705dbec}\\
{\footnotesize\loc{logs/v6/flat_D12.json}} & \texttt{\scriptsize 80804ff73c1b02d05f0b720df14b957d}\newline\texttt{\scriptsize fb73a38e0a9adeaa717e532bc1fcdf84}\\
{\footnotesize\loc{logs/v6/flat_D13.json}} & \texttt{\scriptsize a8008a4a8ac47c9fac9c39efa5729dd7}\newline\texttt{\scriptsize 4fa9d7e70aae824e78407def473578d1}\\
{\footnotesize\loc{logs/v6/flat_vs_p18.json}} & \texttt{\scriptsize 891bf4a356e9427dea75d403874ec422}\newline\texttt{\scriptsize 0d4fa5ba74de59b86ee8152bb010c42c}\\
{\footnotesize\loc{logs/v6/hh_structures.json}} & \texttt{\scriptsize 78fb8461a611975d88dc600ea2f6e937}\newline\texttt{\scriptsize b52a31e62b6ed0483fbf721458b5086a}\\
{\footnotesize\loc{logs/v6/hi_d01.json}} & \texttt{\scriptsize 1ee5a24338961e2fb35006fb00ebee02}\newline\texttt{\scriptsize 0ad93ab20bd78c8f5c0dbe2b7d3dbbb1}\\
{\footnotesize\loc{logs/v6/hi_d01_exact.json}} & \texttt{\scriptsize 850e69e35f7d33cf0933886869cc65f5}\newline\texttt{\scriptsize cd1fdfff1970d8fe1e2bba8ea7979305}\\
{\footnotesize\loc{logs/v6/st_norm.json}} & \texttt{\scriptsize 5b30e5af54e2a46f48f45356d2dee543}\newline\texttt{\scriptsize 4bc378fd19c3a1b4ee5e98fb5f197253}\\
{\footnotesize\loc{logs/v6/st_v6_CR.json}} & \texttt{\scriptsize 9f1065ca1c9ae07f4eb640061540e50e}\newline\texttt{\scriptsize cad8faa0ebc75afd097d5fd943cf8dda}\\
{\footnotesize\loc{logs/v6/st_v6_F.json}} & \texttt{\scriptsize eb09490b5b839a503d9a053fd5d5180a}\newline\texttt{\scriptsize 3473c4507896ef01dd79a7322b23be22}\\
{\footnotesize\loc{logs/v6/st_v6_FLAT.json}} & \texttt{\scriptsize 46b533ea296208474c3893d989fd6c45}\newline\texttt{\scriptsize b3a874acd4e15a92b49e29fc4896bacb}\\
{\footnotesize\loc{logs/v6/st_v6_UNIT.json}} & \texttt{\scriptsize 20a9d7d327e92d8f182e54954aa7bd1d}\newline\texttt{\scriptsize 29cc8995bc35a9803ac239401001be07}\\
{\footnotesize\loc{logs/v6/st_v6_stats.json}} & \texttt{\scriptsize bde0de688cc0e7317f3d42a2cf049758}\newline\texttt{\scriptsize 59d747c16e00c0cfe66cc4bda2ca7711}\\
{\footnotesize\loc{logs/v6/weighted_apex.json}} & \texttt{\scriptsize b23bf74962e3174c689db732c5a4603c}\newline\texttt{\scriptsize 2d380e2be75c58a324dd6a337ead0c2d}\\
{\footnotesize\loc{logs/v6/weighted_checks.json}} & \texttt{\scriptsize bc1b45f39bc41eadb640478f406b2655}\newline\texttt{\scriptsize 51d63dea14a53017d5350d8a582fced5}\\
\multicolumn{2}{@{}l}{\textit{Read-only frozen inputs}}\\
{\footnotesize\loc{RO:XVI-P/code/p5_entry/amplitude_rules.py}} & \texttt{\scriptsize 0daf63a3c8803d54514cb2aac29bea12}\newline\texttt{\scriptsize 7c379e025f32c5b0219d365de5bdde4a}\\
{\footnotesize\loc{RO:XVI-P/code/p5_entry/feynman.py}} & \texttt{\scriptsize 4fab8b0d4a1cc0b4cd0be19ec30ace80}\newline\texttt{\scriptsize 559b92fa9808f1fc25c858599a389344}\\
{\footnotesize\loc{RO:XVI-P/code/p5_entry/dred.py}} & \texttt{\scriptsize 24112d9a6584613474e96f538840a9ff}\newline\texttt{\scriptsize 6aecbf9ebd588cd890fd38d2e295de5a}\\
{\footnotesize\loc{RO:XVI-Q/provenance/XVI-P/certificates/P18_D11.json}} & \texttt{\scriptsize af696a2ad7ec85fbba966e5403079c65}\newline\texttt{\scriptsize f5635d704ec4814f7755ff256d2031e1}\\
{\footnotesize\loc{xviq_regen/Q02_bulk_graphs.json}} & \texttt{\scriptsize 4ac6dc39f6e1763f1a90bceb45296fab}\newline\texttt{\scriptsize d539e086c160d3ea60427350cf118505}\\
\end{longtable}
\endgroup

\paragraph{Rerun commands.}
All commands run from \loc{/workspace/XVI-O1-CONE}.
The spectral runs were launched as detached jobs with logs because
the machine has about 7\,GB of memory.
\begin{center}\footnotesize
\begin{tabular}{@{}l@{}}
\verb!for g in D01 D02 D03 D04 D05a D06 D11 D12 D13; do for b in CR F; do!\\
\verb!  setsid python3 code/spectral_apex.py $g $b > logs/v4/spec_${g}_${b}.log 2>&1; done; done!\\
\verb!for g in D01 D02 D03 D06 D11 D12 D13; do setsid python3 code/spectral_apex.py $g UNIT \!\\
\verb!  > logs/v4/spec_${g}_UNIT.log 2>&1; done!\\
\verb!python3 code/catalogue_check.py          # -> logs/v6/catalogue_check.json!\\
\verb!python3 code/st_norm.py                  # -> logs/v6/st_norm.json!\\
\verb!for b in CR F UNIT; do python3 code/st_v6.py "logs/v4/spec_{g}_$b.json" $b; done!\\
\verb!for g in D01 D02 D03 D06 D11 D12 D13; do python3 code/flat_pole.py $g; done!\\
\verb!python3 code/st_v6.py "logs/v6/flat_{g}.json" FLAT!\\
\verb!python3 code/st_v6_stats.py; python3 code/flat_vs_p18.py!\\
\verb!python3 code/weighted_apex.py; python3 code/weighted_checks.py!\\
\verb!python3 code/hi_d01.py; python3 code/hi_d01_exact.py!\\
\verb!python3 code/hh_structures.py!\\
\verb!python3 code/build_v6.py && python3 code/artifacts_v6.py!\\
\verb!cd v6build && pdflatex XVI-O1-CONE-v6.tex && pdflatex XVI-O1-CONE-v6.tex!
\end{tabular}
\end{center}
The spectral runs (first four lines) were executed during the v4 session
and were not repeated for v6; every other command was run for v6.

%======================================================================
\noindent\emph{Marks:} \(^{*}\) read in full for this revision
(for programme records: the specific files named);
\(^{\ddagger}\) bibliographic details confirmed online;
\(^{\dagger}\) cited from memory or from earlier supplied records,
not re-read here.
\medskip
\begin{thebibliography}{99}
\addcontentsline{toc}{section}{References}
%======================================================================

\bibitem{CheungConeV4}
\(^{*}\)\ N.~Cheung,
\emph{Scoped O.1 on the Cone, v4:
Spectral Apex Reduction, Exact Constant-Profile Apex Poles of All
Two-Point Graphs, the Ghost-Equation Row, and an Unvalidated
\(hh\) Slavnov--Taylor Test},
18 October 2026.
Supplied complete source \loc{XVI-O1-CONE-v4.tex}.
Read in full. Its code and logs are now archived and audited
(Appendix~\ref{app:artifacts}).

\bibitem{CheungConeV5}
\(^{*}\)\ N.~Cheung,
\emph{Scoped O.1 on the Cone, v5: Spectral Apex Reduction,
Constant-Profile Pole Records, and Ward-Identity Boundaries},
18 October 2026. Source \loc{XVI-O1-CONE-v5.tex}
(SHA-256 \loc{e00e49f5...a6985b}, full digest in
Appendix~\ref{app:artifacts}), read in full; the base of this revision.

\bibitem{CheungP}
\(^{*}\)\ N.~Cheung,
\emph{XVI-P},
frozen p22 contract as described in the supplied cone source;
separately supplied p23 excerpt retaining the p21 absolute-D11
normalization and scope discussion.
Enter Calculus Companion, unpublished manuscript records.
The later excerpt is a comparison, not a silent contract update.
Files read in full for v6 (read-only copies, digests equal to the
originals): \loc{code/p5_entry/amplitude_rules.py},
\loc{feynman.py}, \loc{dred.py}.
Relevant scope locators include
\loc{rem:o1}, \loc{rem:p5}, and \loc{rem:qds}.

\bibitem{CheungQ}
\(^{*}\)\ N.~Cheung,
\emph{XVI-Q},
supplied Q02 normalization description and separately scoped
Q11 release record.
Enter Calculus Companion, unpublished source records.
v6 reads the regenerated Q02 graph record
\loc{xviq_regen/Q02_bulk_graphs.json} in full and the
read-only certificate \loc{provenance/XVI-P/certificates/P18_D11.json};
the Q11 scope is retained as supplied.

\bibitem{CheungR}
\(^{\dagger}\)\ N.~Cheung,
\emph{XVI-R},
frozen read-only comparison record identified in
\loc{XVI-O1-CONE-v4.tex}.
Not read for v6; only its full tree digest is recorded
(unchanged before/after).

\bibitem{CheungIB}
\(^{\dagger}\)\ N.~Cheung,
\emph{I-B: A Finite Enter Spacetime:
Calibrated Destructive Time Readout on a Prescribed Conical Product},
version 1.
Supplied version-qualified excerpts, Enter Calculus Companion.
Relevant comparisons:
\loc{def:experiment}, \loc{thm:instrument},
\loc{prop:interfaces}, and \loc{prop:not-source}.
\url{https://enter.grok.me/browse/I-B}.

\bibitem{CheungEn09}
\(^{\dagger}\)\ N.~Cheung,
\emph{En09-A: Realization of Relational Time by Destructive Enter
on a Clock Register},
version 1.
Supplied excerpts, Enter Calculus Companion.
Complete finite record--residual comparison only.
\url{https://enter.grok.me/browse/En09-A}.

\bibitem{CheungIBG}
\(^{\dagger}\)\ N.~Cheung,
\emph{I-B-G: Intrinsic and holonomy-twisted Green operators on a
prescribed conical spacetime:
Self-adjoint extensions, diffraction, finite propagation,
and supported transfer from flat-completion laboratories},
version 2.
Supplied source descriptions and Green-domain excerpts.
Comparisons include \loc{thm:green}, \loc{eq:Carslaw},
\loc{thm:transfer}, and \loc{rem:compact-exact}.
\url{https://enter.grok.me/browse/I-B-G}.

\bibitem{CheungIBGS}
\(^{\dagger}\)\ N.~Cheung,
\emph{I-B-G-S: Intrinsic spin-1 and spin-2 Green operators on the
prescribed conical spacetime},
version 2.
Supplied unpublished excerpts and version-qualified descriptions,
including \loc{thm:spin1}, \loc{thm:spin2},
\loc{def:geo}, and \loc{thm:selection}.
Free source-test compatibility is not interacting loop closure.

\bibitem{CheungTensor}
\(^{\dagger}\)\ N.~Cheung,
\emph{I-B-G-S-B}.
Supplied programme description of positive tensor-core selection,
exceptional source-test defects, matching, and analytical kernel
enclosures.
Comparison locators:
\loc{prop:warddefect}, \loc{cor:exclude},
\loc{thm:contrast}, \loc{lem:modetail},
\loc{lem:complexbound}, and \loc{thm:Cauchy}.
A complete version-qualified proof source was not supplied for
independent audit.

\bibitem{CheungEcho}
\(^{\dagger}\)\ N.~Cheung,
\emph{I-B-ECHO-v2: Diffracted echoes and finite supplied-report
instruments},
version 2, October 2026.
Supplied full-source record and source-qualified comparisons.
Relevant locators:
\loc{lem:front}, \loc{thm:gtd}, \loc{thm:readout},
\loc{thm:TV}, \loc{thm:core},
and \loc{cor:nonreconstruction}.
No continuum acquisition or loop-residue theorem is imported.

\bibitem{CheungVP}
\(^{\dagger}\)\ N.~Cheung,
\emph{I-B-VP: Vacuum polarization of the Friedrichs ground state
on a prescribed conical spacetime},
version 3 and supplied version-qualified predecessor passages.
Enter Calculus Companion.
Comparisons include \loc{lem:infrared} and
\loc{thm:friedrichs}.
One-point state control is not the present loop UV/IR analysis.

\bibitem{CheungVP1}
\(^{\dagger}\)\ N.~Cheung,
\emph{I-B-VP-1: Prescribed apex finite parts on a conical Enter
spacetime: Wick square, stress tensor, and charged current:
strong extension, Ward identities, and classified contact freedom},
version 3, October 2026.
Supplied excerpts and source descriptions.
One-point distributional extension freedom is not a replacement
for a spin-2 operator-domain selection.

\bibitem{CheungEx1B}
\(^{\dagger}\)\ N.~Cheung,
\emph{I-B-VP-5-Ex1-B},
version 1, 18 October 2026.
Supplied normalization, finite-part non-universality, and
verification-scope excerpts.
Relevant locators include
\loc{thm:running}, \loc{thm:continuity},
\loc{thm:equalradius}, and \loc{thm:trace}.
The historical v4 intake is not independently re-executed here.

\bibitem{CheungEx2}
\(^{\dagger}\)\ N.~Cheung,
\emph{I-B-VP-5-Ex2:
Smooth flux-core selection on a prescribed cone:
Scalar Friedrichs universality, Dirac-window selection,
and quantitative exterior matching},
version 1, 18 October 2026.
Supplied source excerpts and bibliography.
Operator selection and Dirac-square comparisons only.

\bibitem{CheungEx2B}
\(^{\dagger}\)\ N.~Cheung,
\emph{I-B-VP-5-Ex2-B},
version 1, 18 October 2026.
Supplied excerpts concerning scalar Mosco selection,
exhausted-annulus arguments, endpoint continuity,
and finite-core qualifications.
No spin-2 CR selection or loop Ward theorem is imported.

\bibitem{CheungXJ}
\(^{\dagger}\)\ N.~Cheung,
\emph{X-J-ST},
version 1.
Supplied unpublished signal and clock instrument excerpts,
including \loc{thm:signal}, \loc{thm:clock},
and \loc{thm:TV}.
Supplied-value reporting and complete failures are retained.

\bibitem{CheungXK}
\(^{\dagger}\)\ N.~Cheung,
\emph{X-K-ST},
version 1.
Supplied unpublished clock-composition and resource-fit excerpts,
including \loc{def:grade} and \loc{thm:release}.
A finite typed tree is not automatically a realized laboratory
or a continuum acquisition process.

\bibitem{CheungIC}
\(^{\dagger}\)\ N.~Cheung,
\emph{I-C-v2: The Shredder Interface:
Destructive Readout, Black-Hole Analogues, and Quantum Computers},
October 2026.
Supplied source in the revision material.
Theorems A--F distinguish referenced recovery, append-only
entropy, conditional thermal balance, classical exterior access,
endpoint underdetermination, and accessible classical persistence.
Documentary separation, not an O.1 premise.

\bibitem{CheungCharacteristic}
\(^{\dagger}\)\ N.~Cheung,
\emph{XIV-H}, version 1;
\emph{XIV-I}, version 1; and
\emph{XIV-J}, version 1.
Supplied excerpts.
Historical H2 ending, constructed I2 internal amplitude,
unresolved I3/J2 full characteristic map, and no inherited
quantum admission.
A partial off-characteristic reference is not the complete map.

\bibitem{CheungXVB}
\(^{\dagger}\)\ N.~Cheung,
\emph{XV-B},
version 2.
Supplied technical-interface excerpts:
reconstruction descent, relative lifting, normal-jet null tests,
and conditional extension.
The technical manuscript is distinguished from the unopened
downstream quantum-admission exit.

\bibitem{Fursaev1994}
\(^{\ddagger}\)\ D.~V.~Fursaev,
``The heat kernel expansion on a cone and quantum fields near
cosmic strings,''
\emph{Classical and Quantum Gravity}
\textbf{11} (1994), 1431--1443.
arXiv:\nolinkurl{hep-th/9309050}.
Apex-local heat coefficients and product-cone context.

\bibitem{FursaevMiele1997}
\(^{\ddagger}\)\ D.~V.~Fursaev and G.~Miele,
``Cones, spins and heat kernels,''
\emph{Nuclear Physics B}
\textbf{484} (1997), 697--723.
DOI: \nolinkurl{10.1016/S0550-3213(96)00631-1}.
arXiv:\nolinkurl{hep-th/9605153}.
Spin-dependent conical heat coefficients and nonuniform smooth limits;
not a substitute for the present graph-specific domain contract.

\bibitem{Vassilevich}
\(^{\ddagger}\)\ D.~V.~Vassilevich,
``Heat kernel expansion: user's manual,''
\emph{Physics Reports}
\textbf{388} (2003), 279--360.
DOI: \nolinkurl{10.1016/j.physrep.2003.09.002}.
arXiv:\nolinkurl{hep-th/0306138}.
Heat-kernel normalization and Laplace-type coefficient background.

\bibitem{DLMF}
\(^{\ddagger}\)\ F.~W.~J.~Olver et al., editors,
\emph{NIST Digital Library of Mathematical Functions},
Chapter 10, especially Sections 10.29, 10.32, and 10.43
(10.43.22 confirmed online for v6).
Bessel recurrences, modified-Bessel relations, and integral transforms.
\url{https://dlmf.nist.gov/10}.

\bibitem{DerezinskiRichard}
\(^{\ddagger}\)\ J.~Derezi\'nski and S.~Richard,
``On Schr\"odinger operators with inverse square potentials
on the half-line,''
\emph{Annales Henri Poincar\'e}
\textbf{18} (2017), 869--928.
DOI: \nolinkurl{10.1007/s00023-016-0520-7}.
arXiv:\nolinkurl{1604.03340}.
Homogeneous Bessel realizations, spectral structure,
and self-adjoint endpoint domains.

\bibitem{KayStuder}
\(^{\dagger}\)\ B.~S.~Kay and U.~M.~Studer,
``Boundary conditions for quantum mechanics on cones and fields
around cosmic strings,''
\emph{Communications in Mathematical Physics}
\textbf{139} (1991), 103--139.
DOI: \nolinkurl{10.1007/BF02102731}.

\bibitem{Stockinger}
\(^{\ddagger}\)\ D.~St\"ockinger,
``Regularization by dimensional reduction:
Consistency, quantum action principle, and supersymmetry,''
\emph{Journal of High Energy Physics}
\textbf{2005}(03) (2005), 076.
DOI: \nolinkurl{10.1088/1126-6708/2005/03/076}.
arXiv:\nolinkurl{hep-ph/0503129}.
Regulator consistency is distinguished from automatic preservation
of a particular identity.

\bibitem{HSR}
\(^{\ddagger}\)\ M.~Horodecki, P.~W.~Shor, and M.~B.~Ruskai,
``Entanglement breaking channels,''
\emph{Reviews in Mathematical Physics}
\textbf{15} (2003), 629--641.
DOI: \nolinkurl{10.1142/S0129055X03001709}.
Preprint title:
``General Entanglement Breaking Channels,''
arXiv:\nolinkurl{quant-ph/0302031}.

\bibitem{NielsenAverage}
\(^{\ddagger}\)\ M.~A.~Nielsen,
``A simple formula for the average gate fidelity of a quantum
dynamical operation,''
\emph{Physics Letters A}
\textbf{303} (2002), 249--252.
DOI: \nolinkurl{10.1016/S0375-9601(02)01272-0}.
arXiv:\nolinkurl{quant-ph/0205035}.
The average-fidelity/entanglement-fidelity relation is attributed
there to earlier Horodecki work.

\bibitem{ReebWolf}
\(^{\ddagger}\)\ D.~Reeb and M.~M.~Wolf,
``An improved Landauer principle with finite-size corrections,''
\emph{New Journal of Physics}
\textbf{16} (2014), 103011.
DOI: \nolinkurl{10.1088/1367-2630/16/10/103011}.
arXiv:\nolinkurl{1306.4352}.
Conditional finite-reservoir entropy balance, not a heat cost
deduced from the abstract readout channel.

\bibitem{Kabat}
\(^{\ddagger}\)\ D.~Kabat,
``Black hole entropy and entropy of entanglement,''
\emph{Nuclear Physics B} \textbf{453} (1995), 281--299.
DOI: \nolinkurl{10.1016/0550-3213(95)00443-V}.
arXiv:\nolinkurl{hep-th/9503016}.
Spin-one horizon contact term; structural analogue only.

\end{thebibliography}

\end{document}
